% ============================================================
% 第3章 提示词工程与智能体系统(扩充版，目标~50页)
% ============================================================

\chapter{提示词工程与智能体系统}


\section*{}
% \addcontentsline{toc}{section}{本章导览}
上一章我们了解了AI如何学习和理解世界，从机器学习到深度学习，从卷积网络到注意力机制。但这些能力要真正为我们所用，离不开人与AI的有效对话。

学期末，小林面临一门课程论文的截止压力。他需要完成一篇关于深度学习在医学影像中的应用的4000字综述。第一次，他打开AI对话框输入了“帮我写一篇论文”，AI只回复了1200字的泛泛之谈。第二次他限定了字数，但引用的都是2020年之前的文献。直到第三次，他仔细拆解了任务，设定了角色、分步给出指令、约束了输出格式，AI终于产出了一篇结构清晰、引用规范的完整综述。同样的AI、同样的话题，只是提问方式不同，结果天差地别。
读完本章，你将找到这些问题的答案。在开始之前，先带着以下几个问题来阅读：

\begin{itemize}
\item \textbf{为什么同一个AI、同样的任务，不同的提问方式会得到天差地别的结果?} 提示词的词法与语法决定了AI理解任务的起点。
\item \textbf{如何让AI按照特定的角色和思维框架来回答问题?} 角色设定不是扮演游戏，而是在AI的知识空间中划定注意力边界。
\item \textbf{AI为什么会忘记前面说过的话?如何管理它的记忆?} 上下文窗口是AI的短期记忆，善用它才能保证连贯性。
\item \textbf{怎样让AI在复杂推理中少犯错误?} 思维链让AI把心算变成笔算，每一步透明可见。
\item \textbf{能否把调教好的提示词保存下来反复使用?} Skill蒸馏将零散的成功经验固化为可复用的能力模块。
\item \textbf{AI如何从你问它答升级为自主完成任务?} 智能体通过感知、规划、工具调用和反馈调整构建执行闭环。
\item \textbf{当多个AI分工协作时，它们如何沟通、如何避免冲突?} 多智能体系统需要主从、议会或市场等组织模式来协调。
\item \textbf{当AI走出屏幕进入物理世界，会面临哪些全新的安全挑战?} 具身智能带来的实时性、隐私权限和责任归属问题。
\end{itemize}

\begin{center}
\textcolor{gray!70}{\*\ \ \*\ \ \*}
\end{center}
\vspace{12pt}

\newpage
% ============================================================
\section[提示词的文法与设计]{\textbf{提示词}\footnote{提示词(Prompt)：用户向AI输入的自然语言指令，引导AI生成特定类型的输出。}的文法与设计}
% ============================================================

\begin{sectionkeypoints}
\item 了解提示词的词法和语法基本规范，理解动词、形容词和约束条件对输出类型与风格的决定性作用
\item 掌握角色设定的本质与五大要素，理解完整的身份、经验、风格、框架和约束对输出稳定性的影响
\item 掌握上下文窗口的概念及其管理策略，理解背景注入、示例引导和格式说明三种控制方式
\item 理解思维链推理的核心思想及零样本思维链、思维树、反思链三种变体的适用场景
\item 掌握提示词综合模板的结构，理解角色设定、分步任务、格式要求和约束条件的整合方法
\item 理解结构化输出的重要性，掌握表格、JSON、Mermaid等常见输出格式的约束技巧
\end{sectionkeypoints}



现在，让我们继续回到小林的例子。有了前面的失败经验，他开始有意识地设计提示词的每一步细节。他这样输入：

\begin{userbox}
你是一位医学影像专业的研究生。请帮我完成一份关于“深度学习在肺结节CT影像诊断中的应用”的文献综述。
第一步搜索近五年(2021-2025年)该领域的核心文献；
第二步按以下框架组织内容：常用深度学习模型(CNN、U-Net、Transformer)、公开数据集(LUNA16、LIDC-IDRI)、关键性能指标(敏感度、假阳性率)、当前局限与未来方向；
第三步为每个部分配2-3个代表性研究的引用；
第四步用表格比较不同模型在LUNA16数据集上的表现。
语言风格保持学术性但不晦涩。
\end{userbox}

\begin{aibox}
\textbf{文献综述：深度学习在肺结节CT诊断中的应用}

\textbf{1. 常用深度学习模型}

CNN类模型(ResNet、DenseNet等)是肺结节检测的基础架构，在LUNA16上取得了92.3\%的敏感度。U-Net及其变体在分割任务中表现突出，Dice系数可达0.89。Vision Transformer在大规模数据上表现优于传统CNN。

\textbf{2. 性能对比}

\begin{formalTable}[H]
\caption{不同模型在LUNA16数据集上的性能对比}
\label{tab:ch3_01_model_performance}
\begin{tabular}{c c c}
\toprule
模型 & 敏感度 & 假阳性率 \\ \midrule
3D CNN & 94.2\% & 1.0/scan \\
U-Net+ResNet & 95.8\% & 0.8/scan \\
Swin Transformer & 96.5\% & 0.6/scan \\ \bottomrule
\end{tabular}
\end{formalTable}

\textbf{3. 当前局限：}小样本泛化能力不足，临床验证研究有限。
\end{aibox}

这一次AI产出了一篇4000多字、结构清晰、引用了十五篇文献的综述。小林拿到AI初稿后，又花了一小时用“临床验证方面能不能增加一些多中心研究的数据”进行第二轮追问，AI补充了三篇2024年的临床试验文献，最后又用了30分钟调整格式和补充个人见解。小林总共约两小时完成了一篇高质量综述。
由此可以得出第一条原则：提示词不是随便问问，而是一种需要学习的沟通语言。

\begin{funbox}
\textbf{Prompt一词的来历。} 

今天AI领域使用的提示词(Prompt)一词，源头可追溯至拉丁语“promptus”，本意为“使准备好、使就绪”。在计算机科学早期，终端界面上的提示符就被称为prompt：它静静地闪烁在屏幕上，等待用户输入指令。2022 年底，随着 ChatGPT 等大语言模型的发布，这个原本属于命令行界面的词汇被赋予了全新的含义：从“系统等待你输入”变成了“你引导系统输出的”魔法语言。历史就这样完成了对同一个词的双重解读。
\end{funbox}

\subsection[词法语法与基本规范]{\textbf{词法}\footnote{词法(Lexicon)：提示词中词汇的选择和使用，包括动词、形容词等关键词类的运用。}\textbf{语法}\footnote{语法(Syntax)：提示词中句子的组织结构和信息排列顺序。}与基本规范}

小林打开AI对话框，最直接的问题就是：该说什么?怎么说?
他和很多人一样，以为提示词就是随便问问，但事实远非如此。提示词有自己的词法：选择合适的词汇；也有自己的语法：组织句子的结构。理解了这两个层面，就能从随便问问升级到精准指挥。 如图\ref{fig:ch3_3_1}所示，结构化提示词通过明确的角色设定、分步任务和格式约束，相比简单一句话提问能获得结构更清晰、内容更丰富的回答。

\begin{figure}[H]
\centering
\BookFigureImage{./images_chapter_01-03/3.1.png}
\caption{一句话提示词与结构化提示词对比}
\label{fig:ch3_3_1}
\end{figure}



\subsubsection*{\textbf{•词法的核心要素}}

词法：选对词。提示词中最关键的词类就是动词。动词决定了AI将用户当作什么角色来服务，也决定了输出内容的类型。例如，同样是问光合作用的知识，如果用不同动词效果完全不同。对AI提问“请描述一下光合作用”，得到的可能是教科书式描述；用“用表格对比C3和C4植物”则会得到结构清晰的对比表格。如图\ref{fig:ch3_3_2}所示，使用“描述”、“对比”、“分析”等不同动词会引导AI产生截然不同的输出格式和信息密度。

\begin{figure}[H]
\centering
\BookFigureImage{./images_chapter_01-03/3.2.png}
\caption{不同动词的输出效果对比}
\label{fig:ch3_3_2}
\end{figure}

再看一个反面案例：如果像这样提问“看一下这个数据”，AI会输出含糊评论，因为它不知道用户要分析、总结还是确认，如下图所示。\par \begin{userbox}
看一下这个数据。

\end{userbox}

\begin{aibox}
好的让我看看。这份数据记录了三个组别的测量结果看起来分布比较均匀。对照组数据波动较大实验组相对稳定。总体来看实验结果有一定参考价值但需要进一步分析。
\end{aibox}

AI不知道用户要分析、总结还是确认。好的提示词应该明确指定任务，并通过
形容词控制风格和深度。如果详细地让AI展开论述，AI可能会写出一整页；如果让AI简洁地提炼要点，通常输出3-5个点；如果让AI以学术风格回答，则会引用文献保持严谨用语；如果让AI以通俗风格回答，则会用比喻和例子解释复杂概念。约束条件的词同样重要：“近五年的文献”限定了时间范围，“中国学术界的观点”限定了地域，“APA格式”指定了引用规范。 

\subsubsection*{\textbf{•语法的组织结构}}

语法：排好序。提示词不是把信息堆砌在一起就行了，句子的组织方式同样重要。

第一种是总分总结构，先给出整体目标，再展开具体任务，最后总结要求，让AI先理解大局再处理细节，输出更有条理；第二种是递进结构，先提供背景信息，再逐步深入，每一步都建立在前一步的基础上；第三种是对比结构，要求AI同时对比多个对象，AI通常在表格中呈现结果，信息密度高、一目了然。 如图\ref{fig:ch3_3_3}所示，总分总、递进和对比三种语法结构各有特点，合理的结构组织能让AI输出更有条理。

\begin{figure}[H]
\centering
\BookFigureImage{./images_chapter_01-03/3.3.png}
\caption{提示词的三种语法结构对比}
\label{fig:ch3_3_3}
\end{figure}





\subsubsection*{\textbf{•常见错误与防范}}


常见的提示词错误有三种。第一种是过于模糊：只说“帮我写点东西”，AI不知道题材、长度、风格，输出随机且几乎不可用。第二种是信息过载：把所有要求一次性塞进一句话，AI无法区分优先级，输出混乱，应该分段落、分要点组织。第三种是缺少约束：不提限制条件，AI可能用2018年的数据(如果需要2025年的数据)，或者用美国市场的案例(如果需要中国市场的案例)。

在对比实验中，小林给AI完全相同的底层问题，但使用不同质量的提示词。结果显示：结构化提示词获得有效回答的概率比简单提问高出约三到四倍。写好提示词不是锦上添花，而是使用AI的必需技能。
\begin{funbox}
\textbf{提示词的两位老祖宗：SHRDLU与ELIZA。} 

1960年代末，MIT接连诞生了两个对AI对话影响深远的程序。一个是约瑟夫·魏泽鲍姆的ELIZA：它通过简单的模式匹配模拟心理治疗师，当用户说“我觉得很沮丧”时，它会回应“你为什么会觉得沮丧?”。尽管魏泽鲍姆明确声明ELIZA没有任何理解能力，许多用户仍然向它倾吐内心最深处的秘密。另一个是特里·维诺格拉德的SHRDLU：它生活在一个虚拟积木世界中，能理解“把那块绿色的方块放到红色棱锥上面”这样的指令。ELIZA教会了我们：人们容易被“听起来像在理解”的程序所迷惑；SHRDLU则证明了：真正的理解需要明确的指令和可执行的分解。今天写提示词时，应当从SHRDLU那里学会精确分解任务，同时警惕ELIZA效应：不要假设AI真的理解我们，而要精确描述我们想要什么。
\end{funbox}

\subsection[角色设定与语义框架]{\textbf{角色设定}\footnote{角色设定(Role Setting)：在提示词中为AI指定身份、能力和行为方式的技巧。}与语义框架}

小林发现，AI默认的回答风格是百科全书式的，中立、全面、不带倾向。但在很多场景中，需要的不是一个百科全书，而是一个特定角色。 如图\ref{fig:ch3_3_4}所示，同一问题下，设定不同角色(专家、教师、顾问)会使AI的回答在知识侧重和表达风格上产生明显差异。

\begin{figure}[H]
\centering
\BookFigureImage{./images_chapter_01-03/3.4.png}
\caption{不同角色设定的回答差异}
\label{fig:ch3_3_4}
\end{figure}

\subsubsection*{\textbf{•角色设定的本质与要素}}


角色不是AI真的变成了那个人，而是为AI提供了一个思维框架：限制它的知识范围、调整它的表达方式、指定它的判断标准。可以把角色理解为一个知识过滤器：当AI接收到角色指令后，它会在庞大的参数空间中激活与该角色相关的知识子集，同时抑制不相关的知识。

这个机制解释了为什么同一问题设定不同角色会得到截然不同的回答：不是AI在扮演，而是它真的从不同的知识子集中提取信息、用不同的表达方式组织语言。
\begin{funbox}
\textbf{从注意力到角色设定：划定AI的知识边界。} 

理解角色设定为什么要先理解注意力?从第2.4节可知，注意力机制是AI理解信息的核心：它决定模型看哪里。角色设定的本质，正是利用这个机制为AI划定注意力的边界：当告诉AI“你是一位资深内科医生”时，是在指示它的注意力聚焦于医学知识路径，抑制无关信息。当告诉AI“请用SWOT框架分析”时，是在约束它的注意力在优势→劣势→机会→威胁四个维度的分配方式。


\end{funbox}

当告诉AI“你是一位资深面试官”时，模型会激活与面试相关的知识体(常见的提问模式、评价标准、沟通话术)，同时抑制其他不相关的知识。

角色设定的应用实例：
假设我们要让AI完成几项不同的任务。如果告诉AI“你是一位资深面试官，正在招聘一名前端工程师，请根据这份简历列出5个技术面试问题并附上期望回答要点”，AI输出的问题会集中在JavaScript、React、性能优化等技术领域，语气专业且具有甄别性。如果告诉AI“你是一位大学英语四级考官，请生成一篇四级阅读理解短文并附带5道选择题”，AI会自动调整词汇难度至CET-4水平，设计干扰项合理的选项。如果告诉AI“请你扮演一位历史学教授，解释唐朝安史之乱的起因和影响”，AI会以学术口吻回答，引用史料并给出不同史学流派的观点对比。

角色不只是身份标签，更需要配套能力描述、语气风格和约束条件。例如：

\begin{lightquote}
“你是一位有15年经验的儿科医生，擅长用通俗易懂的语言向家长解释病情。在给出诊断建议前，请先列出所有可能性并按概率排序。不要使用专业术语，除非给出对应的通俗解释。”
\end{lightquote}

这个角色设定包含了身份(儿科医生)、经验年限(15年)、沟通风格(通俗易懂)、思考框架(先列可能性再排序)、约束条件(术语需附带解释)。越完整的角色设定，AI的输出越稳定。 如图\ref{fig:ch3_3_5}所示，有效的角色设定包含身份、经验年限、沟通风格、思考框架和约束条件五大要素，要素越完整输出越稳定。

\begin{figure}[H]
\centering
\BookFigureImage{./images_chapter_01-03/3.5.png}
\caption{有效角色设定的五大要素}
\label{fig:ch3_3_5}
\end{figure}

反面案例：角色设定不足。有些角色设定看似给出了身份，实际上信息不足。例如，如果只说“你是一位专家，帮我分析一下这个市场”，AI不知道是哪方面的专家、什么级别的专业度、以什么立场来分析。又如，“请以CEO的口吻回答”只有身份没有配套的能力描述，AI虽然语气像CEO，但分析深度可能不够。


\subsubsection*{\textbf{•语义框架的应用}}
除了角色设定，还可以指定思考框架，要求AI按特定的方法论分析问题。例如SWOT框架要求从优势、劣势、机会、威胁四个维度展开评估；金字塔原理要求先给出核心结论，再逐层展开论据，每个论据配具体数据支撑；5W1H框架则要求从谁、什么、何时、何地、为何、如何六个角度分析问题。


\subsubsection*{\textbf{•多角色切换策略}}
一个强大的技巧是在同一对话中切换角色。先让AI扮演“严苛的审查者”挑出方案的漏洞，再让它扮演“乐观的推行者”提出实施建议。 小杨在准备课题汇报时，先让AI扮演同行评审专家，逐条质疑研究方案的方法论合理性，AI指出了3个潜在的设计缺陷；然后切换角色，让AI扮演跨学科专家，从数据分析角度重新评估方案，AI又补充了两个之前忽略的交互效应因素。通过两个角色的交叉审视，小杨的最终方案既规避了风险，又抓住了机遇。这种多角色分析能让决策更全面，不需要换人，只需要换提示词。

\subsection{上下文控制与作用机制}

小林有时会遇到这样的情况：明明刚才已经和AI说过某件事了，过了一会儿再问，它好像完全忘了。这不是AI的记忆力差，而是因为一个限制：\textbf{上下文窗口}\footnote{上下文窗口(Context Window)：AI模型一次能处理的文本总量上限，超出部分会被截断或遗忘。}。 如图\ref{fig:ch3_3_6}所示，上下文窗口中靠近开头和结尾的信息具有最高的注意力权重，AI对中间部分的记忆相对较弱。

\begin{figure}[H]
\centering
\BookFigureImage{./images_chapter_01-03/3.6.png}
\caption{上下文窗口与注意力优先级}
\label{fig:ch3_3_6}
\end{figure}

\subsubsection*{\textbf{•上下文窗口的概念与演进}}


上下文窗口是AI模型一次能处理的文本总量。以通义千问为例，它的上下文窗口达到百万级(约相当于75万汉字)。如果对话总长度超过这个限制，最早的信息就会被遗忘。


\vfill
\begin{knowledgebox}
上下文窗口这一概念源于Transformer架构的核心机制：自注意力。在谷歌团队2017年发表的论文《Attention Is All You Need》中，Transformer通过注意力机制让每个词都能关注序列中的所有其他词，但其计算复杂度随序列长度平方增长。早期模型的上下文窗口仅有512个token(约400个汉字)，大语言模型扩展到2048个token，通义千问达到了百万级 token，而Gemini 1.5 Pro已支持1M token。每一次窗口扩展，都意味着AI能同时记住的信息量成倍增长。
以下为各主流模型上下文窗口的演进历程：

\begin{formalTable}[H]
\caption{主流大模型上下文窗口演进}
\label{tab:ch3_02_context_window}
\begin{tabular}{c c c c}
\toprule
\bfseries 模型 & \bfseries 发布时间 & \bfseries 上下文窗口 & \bfseries 约当汉字数 \\ \midrule
Transformer(原版) & 2017 & 512 tokens & 400字 \\
大语言模型 & 2020 & 2，048 tokens & 1，600字 \\
通义千问 & 2022 & 4，096 tokens & 3，200字 \\
% GPT-4 & 2023 & 8，192 tokens & 6，400字 \\
% Claude 2 & 2023 & 100K tokens & 8万字 \\
通义千问 & 2023 & 128K tokens & 10万字 \\
% Gemini 1.5 Pro & 2024 & 1M tokens & 80万字 \\ \bottomrule
\end{tabular}
\end{formalTable}

每一次窗口扩展都意味着AI能同时记住更多的信息，从一页纸到一座图书馆的内容。但更大的窗口也带来了更高的计算成本，需要在实际应用中权衡记忆容量与响应速度。
\end{knowledgebox}
\vspace{1em}


\subsubsection*{\textbf{•三种上下文控制方式}}

了解这个限制后，需要学会管理上下文，常见的方法有以下几种。

\textbf{背景注入}：在提示词开头主动提供必要的背景信息。“基于以下背景信息，请回答......”即使前面已经讨论过，也要重新说明，确保AI始终知道来龙去脉。

\textbf{示例引导}：给AI两三个输入输出的示例，它就能模仿示例的格式和逻辑处理后续输入。比如给两个用户评论和对应AI回复的例子，AI就会按照同样的格式回复其他评论。

\textbf{格式说明}：告诉AI输入数据的结构和含义。“以下数据中，第一列是日期，第二列是销售额，第三列是同比变化率。请分析第三季度的销售趋势。”明确的格式说明避免了AI解析错误。 如图\ref{fig:ch3_3_7}所示，背景注入、示例引导和格式说明三种上下文控制方式各有不同的工作机制和适用场景。

\begin{figure}[H]
\centering
\BookFigureImage{./images_chapter_01-03/3.7.png}
\caption{三种上下文控制方式}
\label{fig:ch3_3_7}
\end{figure}

% \textbf{信息优先级。} 研究表明，提示词中越靠近末尾的内容，AI的注意力权重越高。利用这个特性，可以将最重要的约束条件放在提示词的最后：“请根据以上信息撰写一份分析报告。特别需要注意的是：所有数据和建议必须基于2025年之后的信息，无论如何不要引用2020年之前的数据。”


\subsubsection*{\textbf{•上下文管理实用策略}}


上下文管理的实用策略有三种。一是分段对话，将长任务拆分为多次短对话，不要在一次对话中要求AI做十件事，而是分三四次，每次聚焦两三个任务。例如小林需要AI辅助完成一篇文献综述，他没有一次性把所有要求都写进提示词，而是先让AI搜索和整理核心文献，得到结果后再让AI分析文献之间的相互关系，最后才让AI基于前两步的结果撰写综述，这种分段策略使AI在每个阶段都能集中处理有限信息，输出质量明显优于一次性完成。

二是关键信息重复，在每次新对话中都主动重复关键背景信息。例如，小林在使用AI辅助撰写文献综述时，每次开启新对话都会先补充一句“基于我们之前讨论的肺结节CT诊断方向，以下是新的问题”，即使AI在上一轮对话中已经了解这些背景，这种重复能确保AI始终在正确的知识框架内回答问题，避免因上下文窗口溢出而丢失关键信息。

三是使用RAG（Retrieval-Augmented Generation，检索增强生成）技术，对于需要大量背景知识的任务，将文档提前存入知识库，让AI在需要时检索相关内容。例如，一家医院的知识库中存储了数千份病历资料，当医生向AI询问“请根据类似病例给出治疗方案建议”时，AI会从知识库中自动检索最相关的历史病例，而不是仅依赖训练数据中的普适性知识。这种策略既能处理超长文本的限制，又能确保回答基于最新的专业资料。

\subsection{思维链推理与过程设计}

小林发现，AI在面对需要多步推理的问题时，经常犯一个奇怪的错误：它直接跳到最后一步，给出了一个看似合理但实际上是错误的答案。一个经典的例子：
\begin{funbox}
\textbf{让AI一步步思考的神奇咒语。} 

2022年，日本研究者Kojima等人发表了一篇看似不起眼的论文：他们发现只需在提示词末尾加上一句“让我们一步步思考”，就能让AI的复杂推理准确率大幅提升。这一现象暗示了一个深层次的问题：大语言模型其实知道怎么推理，只是需要被引导出来。这个发现后来被称为零样本思维链，它彻底改变了人们使用AI的方式：从给我答案变成了请展示你的思考过程。
\end{funbox}


“Q：一个池塘里的荷花每天翻倍生长，第30天覆盖满池塘。问第几天覆盖一半?A：第15天。”这是一个典型的错误答案--正确答案应是第29天。但如果我们在提示词末尾加上“请一步步推理”，结果就完全不同了：“Q：......请一步步推理。A：第30天覆盖满池塘。因为每天翻倍，所以第29天覆盖一半。第28天覆盖1/4。验证：第29天的一半到第30天翻倍变成满池塘，符合条件。所以答案是第29天。”

如图\ref{fig:ch3_3_8}所示，面对同一推理问题时，直接回答容易给出错误答案，而思维链通过逐步推理能准确找到解答。

\begin{figure}[H]
\centering
\BookFigureImage{./images_chapter_01-03/3.8.png}
\caption{思维链与直接回答效果对比}
\label{fig:ch3_3_8}
\end{figure}

\subsubsection*{\textbf{•思维链的核心思想}}


这就是\textbf{思维链(Chain-of-Thought， CoT)}\footnote{思维链(Chain-of-Thought， CoT)：要求AI展示逐步推理过程的技术，能显著提升复杂问题的准确率。本质上是通过外部化推理过程来减少推理错误。}的核心思想：不直接要答案，而是要求AI展示完整的推理过程。这相当于让AI把心算变成笔算，每一步都写下来，错误更容易被发现和纠正。
思维链为什么有效?可以从其数学本质来理解。设输入为$x$，期望输出为$y$。直接推理可以表示为：

\[
y = f(x)
\]

其中$f$是模型内部的一次性映射。而思维链将这个过程分解为中间推理步骤$z_1, z_2$等：

\[
z_1 = f_1(x),  \quad z_2 = f_2(x, z_1)， \quad \dots， \quad y = f_k(x, z_{k-1})
\]

每个$f_i$解决一个子问题，比直接求解$f$简单得多。研究表明，思维链能将复杂推理任务的准确率提升约30\%\footnote{Wei J, Wang X, Schuurmans D, et al. Chain-of-Thought Prompting Elicits Reasoning in Large Language Models. NeurIPS, 2022.}。如图\ref{fig:ch3_3_9}所示，思维链将输入拆解为多个中间推理步骤，每个步骤解决一个子问题，从而逼近最终答案。

\begin{figure}[H]
\centering
\BookFigureImage{./images_chapter_01-03/3.9.png}
\caption{思维链逐步推理过程}
\label{fig:ch3_3_9}
\end{figure}

\subsubsection*{\textbf{•思维链的三种变体}}

思维链有几种变体。零样本思维链是最简单、最通用的方式，只需在提示词末尾加一句“让我们一步步思考”即可，不需要提供任何示例，AI会自动进入推理模式。

思维树(Tree-of-Thoughts， ToT)则同时探索多个推理分支，评估每个分支的可行性后选择最优路径，适用于方案设计、策略规划等需要多想几条路的任务。

反思链是AI执行完任务后进行自我评估和修正的技术，适用于写作润色、代码调试等需要多次迭代的任务。 


\subsubsection*{\textbf{•适用场景与选择原则}}

思维链最适用于需要多步推理、逻辑推导的任务。例如，对于数学推理问题“一个商品打8折后再打9折相当于打了几折”，通过一步步推理得到0.8乘以0.9等于0.72，即72折，而直接问可能得到错误答案。又如方案比较任务，让AI逐条列出方案A和B的优劣势并给出判断依据，结果的可信度更高。再如代码调试任务，让AI逐行解释代码逻辑并指出可能出错的地方，更容易定位bug。


三种变体的适用场景各有侧重。零样本思维链最简单、最通用，适合数学题、逻辑题、需要因果分析的问题。思维树同时探索多个推理分支，适合方案设计、策略规划等开放性问题。反思链则适合写作润色、代码调试、报告审校等需要写完再检查的任务。

实际使用中，这三种方法可以灵活组合。例如在方案设计中使用思维树生成多个选项，再用反思链评估每个选项的质量。小张在确定课题研究方案时，先使用思维树生成了5个候选方案，然后用反思链逐一评估每个方案的可行性，最终筛选出最优方案。


\subsubsection*{\textbf{•思维链的局限性}}


不是所有问题都需要思维链。对于简单的事实性问题，思维链只会浪费时间和计算资源。判断是否需要思维链的一个简单标准是：如果这个问题可以用一步推理完成，就不需要思维链；如果需要两个以上的逻辑步骤，思维链会有帮助。

\begin{funbox}
\textbf{高斯巧算：跳步的代价。} 

数学天才高斯10岁时，老师给全班出了一道题：从1加到100等于多少?其他孩子埋头逐项加算，高斯几秒后写出了5050。他发现的规律是：(1+100)×50=5050。这个巧算如此精妙，以至于老师都惊叹不已。但如果让AI来解这道题，不给思维链约束，它可能直接跳步到5050，结果虽然正确，但完全看不到推理过程。读者无法验证这个结果是否来自严谨的数学推导，还是来自蒙对了。高斯的巧算恰恰是思维链的反面教材：他的解法虽然优美，但直接跳过了中间步骤，掩盖了思考过程。而AI的思维链要求恰恰相反：不要跳步，每一步都要写下来，让推理过程透明可见。在提示词工程中，这一原则同样适用于所有复杂任务：宁可繁琐一点，也要让AI展示完整的思考过程。
\end{funbox}

\subsection[结构化输出与格式规范]{\textbf{结构化输出}\footnote{结构化输出(Structured Output)：要求AI按特定格式(表格、JSON、Mermaid等)输出内容的技巧。}与格式规范}

小林好不容易让AI给出了高质量的回答，但打开一看，一大段文字混在一起，关键数据埋没在段落里，费了好大劲才把需要的信息提取出来。
问题的根因在于：没有告诉AI输出应该长什么样。 

\subsubsection*{\textbf{•结构化输出的意义}}


为什么结构化输出很重要?AI的原始输出是自然语言，人类虽能读懂，但机器解析困难。指定输出格式后，AI的结果可以直接被程序消费，如存入数据库、导入电子表格、嵌入网页。在需要批量处理大量任务的场景中，结构化输出不是锦上添花，而是必须项。如图\ref{fig:ch3_3_10}所示，同一份数据以无格式段落呈现时，关键信息埋没在文本中，以结构化表格呈现则一目了然。

\begin{figure}[H]
\centering
\BookFigureImage{./images_chapter_01-03/3.10.png}
\caption{无格式与结构化输出对比}
\label{fig:ch3_3_10}
\end{figure}

\subsubsection*{\textbf{•常见结构化格式}}
常见的结构化格式包括以下几种：Markdown表格适合对比展示，例如可以这样和AI对话：“请用Markdown表格比较A、B、C三个方案的成本、周期和风险”，输出可直接用在文档和网页中。

JSON适合程序解析，例如“请用JSON格式返回”，输出可被程序直接解析；列表/编号适合步骤说明，例如“请用编号列表列出实施这个方案所需的5个步骤”，清晰易读。

Mermaid图表适合可视化，例如“请用Mermaid语法画出这个业务流程的流程图”，直接生成可渲染的图表代码。 如图\ref{fig:ch3_3_11}所示，指定表格、JSON、Mermaid等结构化格式后，AI的输出可以直接被程序消费和使用。

\begin{figure}[H]
\centering
\BookFigureImage{./images_chapter_01-03/3.11.png}
\caption{结构化输出效果对比}
\label{fig:ch3_3_11}
\end{figure}


\subsubsection*{\textbf{•格式约束技巧}}


格式约束的技巧有三种。其一，指定分隔符，用特定的符号分隔不同部分，明确结构边界。其二，给出模板，给AI一个具体的空白模板，它会严格参照模板的结构填充内容。其三，指定长度，例如“每条建议不超过50字”或“正文总字数控制在800字以内”。

反面案例：小杨是一名研一的学生，他需要整理50篇相关文献的摘要用于文献综述。第一次他直接用“帮我总结这篇论文的核心观点”，结果50篇摘要风格不统一：有的偏重研究方法，有的偏重结论应用，有的写了100字，有的却写了400字。他花了两小时重新整理才勉强统一了格式。第二次他吸取教训，设计了通用的模板：“论文标题、研究问题、方法概述（1\~{}2句）、主要发现（2\~{}3句）、结论与意义（1\~{}2句）、关键词（3\~{}5个）。请用客观学术语言，每篇摘要控制在200字以内。”AI严格按模板输出，50份文献摘要结构一致，直接用于文献综述，前者的时间成本是后者的数倍。

% 学好了一对一的提示词技巧后，下一个问题是：每次做相似的任务都要重写一遍提示词吗?

% 小杨在使用AI完成一次实验报告的数据分析后，要求AI以JSON格式输出核心指标，包括均值、标准差、样本量、p值和置信区间。AI生成的JSON可以直接导入数据分析软件进行后续处理，而无需手动转录数据。这让他意识到：当输出格式为机器可读时，AI生成的数据就不再是一次性的，而可以成为整个工作流中的一个标准环节。有没有办法把打磨好的提示词保存下来，下次直接复用?这就引出了下一节：Skill模块。

% ============================================================

\subsection*{提示词综合模板}

至此学习了提示词的各个要素：词法选择、语法结构、角色设定、上下文控制、思维链推理和结构化输出。将这些要素整合在一起就形成了一份通用的提示词模板：

\subsubsection*{\textbf{•通用提示词模板}}


\begin{userbox}
【角色设定】你是一位[具体角色]有[经验年限]年[领域]经验。
【背景信息】[提供任务相关的上下文和数据]
【具体任务】
第一步[子任务1的描述]；
第二步[子任务2的描述]；
第三步[子任务3的描述]；
【输出要求】请用[格式要求]输出[风格要求]。
【约束条件】注意：[最重要的限制条件]。
请一步步推理。
\end{userbox}


\subsubsection*{\textbf{•模板应用案例}}
\textbf{模板应用案例：}

% \footnote{Kim et al.， Nature Biotechnology， 2022.}\footnote{Lander， New England Journal of Medicine， 2021.}

\begin{userbox}
你是一位生物伦理学方向的研究助理，有3年基因编辑伦理研究经验。
我正在进行一项关于CRISPR技术在人类胚胎编辑中的伦理争议的研究。
请帮我完成以下任务：
第一步搜索近五年(2021-2025年)关于CRISPR人类胚胎编辑的核心伦理争论；
第二步将争论按安全性风险、基因多样性、社会公平、监管框架四个主题分类；
第三步为每个主题总结2-3个核心观点和代表性专家的立场；
第四步用表格呈现各方观点的对比。
请用学术论文的语言风格输出参考文献按APA格式列出。
注意：仅纳入经同行评审的学术期刊论文不要引用预印本或媒体报道。
\end{userbox}

\begin{aibox}
\textbf{CRISPR技术在人类胚胎编辑中的伦理争议}

\textbf{1. 安全性风险}
支持方认为随着技术进步脱靶效应已大幅降低。反对方指出对人类胚胎的编辑可能产生不可预知的跨代影响目前对胚胎发育的理解仍不完整。

\textbf{2. 基因多样性}
部分学者担忧基因编辑可能导致人类基因库的单一化减少自然多样性。也有学者认为这只限于极少数个体的编辑不会对整个人类基因库产生显著影响。

\textbf{3. 各方观点对比}
\begin{formalTable}[H]
\caption{基因编辑各方观点对比}
\label{tab:ch3_03_debate}
\begin{tabular}{c c c}
\toprule
主题 & 支持方观点 & 反对方观点 \\ \midrule
安全性 & 技术已足够成熟 & 跨代影响未知 \\
基因多样性 & 影响微乎其微 & 长期风险不可控 \\
社会公平 & 可通过政策调节 & 必然加剧不平等 \\ \bottomrule
\end{tabular}
\end{formalTable}

% \textbf{参考文献}：Kim et al. (2022) Nature Biotechnology； Lander (2021) New England Journal of Medicine.
\end{aibox}

这个模板把前面学到的各个技巧融合在了一起：角色设定让AI进入专业状态、分步任务让AI按逻辑推进、格式要求让输出可读性高、约束条件过滤掉不相关信息。在实践中可以根据具体需求调整每个部分的内容。

\begin{center}
\textcolor{gray!70}{\*\ \ \*\ \ \*}
\end{center}
\vspace{12pt}



\subsubsection*{小结}

提示词工程是从“随便问问”到“精准指挥”的核心技能。掌握词法选择、语法组织、角色设定、上下文控制、思维链推理和结构化输出等六大核心要素，可以显著提升AI的输出质量和稳定性。从一句话提问到结构化提示词，每一次优化都让AI更好地理解用户意图，产出更精准可靠的内容。

\bigskip
\subsection*{课后习题}

\begin{enumerate}
\setlength{\itemsep}{6pt}
\item 提示词工程中词法和语法分别指什么？各举一个应用实例。
\item 角色设定的本质是什么？有效的角色设定应包含哪五大要素？
\item 什么是上下文窗口？它的大小对AI对话有什么影响？
\item 思维链（CoT）推理为什么能提升复杂问题的准确率？它的数学本质是什么？
\item 结构化输出有哪些常见的格式？分别适用于什么场景？
\end{enumerate}

\newpage

\section[Skill模块的抽取复用]{\textbf{Skill}\footnote{Skill(技能模块)：封装了提示词、工具和流程的可复用能力单元。}模块的抽取复用}
% ============================================================

\begin{sectionkeypoints}
\item 理解Skill是提示词、工具和流程三者的结合体，掌握其与插件、工具和Agent的概念分层
\item 掌握Skill蒸馏的三步法(回顾、归纳、泛化)，能从日常AI交互中提炼可复用的模板
\item 了解组合Skill的流水线工作方式，理解多个子Skill串联完成复杂任务的协调机制
\item 理解参数化调用的价值，掌握同一Skill通过配置不同参数适配多种场景的方法
\item 了解Skill的分类体系、版本管理和分享协作对团队效率的提升作用
\item 熟悉Coze、Dify等平台等主流平台的Skill封装与调用方式
\end{sectionkeypoints}




小林花了一个下午精心打磨了一条用于自动整理课程笔记的提示词，从录音转文字到提取关键词，再到按章节分类，最后生成复习提纲，每一步都考虑到了。室友看到了效果说：“你这个整理笔记的方法能教教我吗?”。室友复制粘贴了提示词，发现非常好用，结果隔壁班的同学知道了也来要。

小林这时候意识到：这条精心打磨的提示词本质上是一个可以反复使用的能力模块。如果能固化下来、绑定工具、设计好输入输出接口它就变成了一个Skill。不需要每次重写只需要一键调用。

\subsection{Skill概念与基础定义}

什么是Skill?Skill不是一条简单的提示词，而是三者的结合体：

% \begin{quote}
\textbf{Skill = 提示词 + 工具 + 流程}
% \end{quote}

\vfill
\begin{knowledgebox}
Skill一词在AI领域的用法，很大程度上借用了游戏中的技能树(Skill Tree)概念：玩家通过解锁技能来获得新的能力，不同技能可以组合使用。2018年，在《王者荣耀》的AI项目多智能体系统中首次大规模使用模块化能力单元来描述AI可以执行的不同任务类型。此后，Coze的插件、Dify的工具、乃至智能体主流框架的Chain等平台，都纷纷采用了相似的封装思想。尽管叫法不同：插件、工具、技能、Skill，本质都是将提示词、工具和流程打包为一个可复用的能力单元。 如图\ref{fig:ch3_3_12}所示，Skill由提示词、工具和流程三大要素组成，三者结合构成一个完整可复用的能力单元。
\end{knowledgebox}
\vspace{1em}



\begin{figure}[H]
\centering
\BookFigureImage{./images_chapter_01-03/3.12.png}
\caption{Skill的三大组成要素}
\label{fig:ch3_3_12}
\end{figure}

\subsubsection*{\textbf{•Skill的三大组成要素}}


一是提示词，经过优化的指令模板，用户只需要输入核心参数。比如一个会议纪要Skill的提示词模板中包含“会议主题：{ }、参与人：{ }、会议日期：{ }”等占位符，用户只需填入具体会议信息即可。

二是工具，Skill需要的外部能力：搜索最新信息、读取数据库、发送邮件、生成图表等。没有工具，AI只能依靠训练数据中的知识；有了工具，AI能获取实时信息和操作外部系统。例如小杨的天气查询Skill需要调用气象API获取实时数据，而文献检索Skill则需要连接PubMed等学术数据库。

三是流程，多个步骤的串联：先做什么，再做什么，什么条件下做什么。流程定义了Skill的执行逻辑，确保输出质量的一致性。 如图\ref{fig:ch3_3_13}所示，使用Skill之前手动完成一项任务需要多步操作，封装为Skill后仅需输入核心参数即可自动完成。

\begin{figure}[H]
\centering
\BookFigureImage{./images_chapter_01-03/3.13.png}
\caption{Skill使用前后效果对比}
\label{fig:ch3_3_13}
\end{figure}


\subsubsection*{\textbf{•竞品分析Skill案例}}

小杨每周都要做研究动态追踪。传统的做法：打开学术数据库、搜索最新论文、阅读3-5篇摘要、整理关键进展、撰写分析、排版。整个过程大约2-3小时。
封装为Skill后：
(1) 用户在Skill界面输入：“帮我分析本周特斯拉和比亚迪的动态”，

(2) Skill自动调用网络搜索工具，获取本周两家公司的最新新闻，

(3) Skill按预设的分析框架(产品发布、市场策略、财务表现、行业评价)整理信息，例如对于“特斯拉本周动态”，Skill会自动抓取新品发布信息并归类到产品发布，从财报中提取销量和利润数据归入财务表现；

(4) Skill将分析结果填充到预设的周报模板中，

(5) 用户收到结构完整的竞品分析周报，全程耗时约5分钟，且格式一致。


\subsubsection*{\textbf{•概念分层辨析}}
概念分层：Skill vs 插件 vs 工具 vs Agent。
\textbf{插件}是以标准化接口接入AI的外部应用。比如天气插件、航班查询插件。插件是Skill可以调用的零件。\textbf{工具}是AI可以直接调用的函数或API，如计算器、搜索引擎、代码解释器，多个工具可以组合使用。\textbf{Skill}是封装了提示词、工具和流程的完整能力单元，它调用工具但本身不是工具，包含了怎么用工具、为什么用、用完后怎么处理结果等更高层次的逻辑。\textbf{Agent(智能体)}是能自主决策调用哪些Skill来完成目标的高级系统，Agent不仅会用工具，也会用Skill。

在现实平台中，Skill有不同的实现形式：
Coze平台的插件允许用户创建可复用的能力模块，每个插件包含触发词和参数定义。Dify平台的工具支持用户在可视化界面中配置提示词模板、关联知识库、设置输出变量。智能体市场允许用户上传指令、知识和能力配置，创建专用AI助手。等主流框架的Chain将多个处理步骤打包为可调用的模块，支持串联和并行执行。

\subsection[Skill蒸馏与抽取方法]{\textbf{Skill蒸馏}\footnote{Skill蒸馏(Skill Distillation)：从AI交互经验中提取共性模式，构建可复用模板的方法。}与抽取方法}

已经用AI完成过很多任务，其中有些效果特别好的对话：AI回复的质量远超预期，几乎可以原样使用。如何把这些零散的成功经验提炼成可复用的Skill?
\textbf{Skill蒸馏}就是把日常工作中效果好的AI交互，提炼为通用模板的方法。分为三步：

\vfill
\begin{knowledgebox}
蒸馏(Distillation)一词在机器学习中有两个层面的含义。2015年，Geoffrey Hinton（杰弗里·辛顿）提出了知识蒸馏(Knowledge Distillation)技术：用一个大型教师模型来教一个小型学生模型，将大模型的知识压缩转移给小模型，就像化学蒸馏中将混合物分离提纯一样。而本章讨论的Skill蒸馏则是在一种隐喻层面借用这一概念：从大量AI交互经验中提取共性模式，剔除具体业务细节后凝练为可复用的模板。两种蒸馏的共同点是：去除杂质，保留精华。 如图\ref{fig:ch3_3_14}所示，知识蒸馏将大模型知识压缩到小模型，而Skill蒸馏从AI对话历史中提取共性模式构建可复用模板，两者都是去除杂质保留精华。
以下对比了两种蒸馏的异同：

\begin{formalTable}[H]
\caption{知识蒸馏与Skill蒸馏对比}
\label{tab:ch3_04_distillation}
\begin{tabular}{c c c}
\toprule
\bfseries 维度 & \bfseries 知识蒸馏 & \bfseries Skill蒸馏 \\ \midrule
\bfseries 目的 & 大模型压缩为小模型 & 交互经验固化为模板 \\
\bfseries 输入 & 教师模型参数 & AI对话历史 \\
\bfseries 输出 & 学生模型参数 & 通用提示词模板 \\
\bfseries 方法 & 软标签蒸馏 & 回顾$\rightarrow$归纳$\rightarrow$泛化 \\
\bfseries 效果 & 性能接近大模型但体积小 & 输出质量接近最优对话 \\ \bottomrule
\end{tabular}
\end{formalTable}
\end{knowledgebox}
\vspace{1em}

\begin{figure}[H]
\centering
\BookFigureImage{./images_chapter_01-03/3.14.png}
\caption{Skill蒸馏流程}
\label{fig:ch3_3_14}
\end{figure}

\subsubsection*{\textbf{•Skill蒸馏的三步法}}


第一步：回顾。翻看历史对话，找出那些效果特别好的交互。判断标准是：AI的回答基本不需要修改就能用，或者稍作修改就能达到专业水平。把这些对话标记出来，作为蒸馏的原料。

第二步：归纳。逐条分析这些成功对话，提取共性：什么角色设定最有效?什么输出格式最常用?什么约束条件最关键?不要关注具体的业务内容(分析iPhone和三星)，而要关注模式(用表格对比两个产品的功能参数)。

第三步：泛化。去掉具体的业务细节，保留可复用的框架。将“告诉我关于iPhone 16的相机参数”泛化为“告诉我关于{产品名}的{属性}参数”。将“分析特斯拉和比亚迪的竞争优势”泛化为“对比分析{公司A}和{公司B}在{行业}中的{竞争维度}。 如图\ref{fig:ch3_3_15}所示，Skill蒸馏经历从具体案例到模式归纳再到通用模板的三阶段转化过程。”

\begin{figure}[H]
\centering
\BookFigureImage{./images_chapter_01-03/3.15.png}
\caption{Skill蒸馏过程：案例$\rightarrow$模式$\rightarrow$模板}
\label{fig:ch3_3_15}
\end{figure}


\subsubsection*{\textbf{•从实例到模板的转化}}

小张做过10次润色学术论文摘要的任务，每次效果都不错。回顾发现，每次都要求了学术风格、摘要结构和专业术语等方面的规范化处理。把这些要求固化为产品润色Skill，以后每次只需要粘贴原始描述文本即可。不需要再思考要提哪些要求：Skill已经封装好了。

处理用户反馈时最好的回复结构是先共情、再解释、然后给出解决方案、最后邀请再次体验。将这个结构固化为反馈回复Skill，所有同学统一使用，不仅效率提升，服务质量也更加稳定。


\subsubsection*{\textbf{•挖掘Skill素材的信号}}
不是所有经验都必须从历史对话中提取，也可以主动发现这个任务并做成一个Skill。下面这些信号说明找到了一个潜在的Skill：

一是这个任务每周都做，频率高，值得封装。例如每周竞品简报，对应竞品动态追踪Skill。
二是这个任务很耗时但模式固定，重复性高，封装后效率提升大。例如每次汇报前的数据可视化，对应数据图表建议Skill。
三是每次做的方法基本相同只是参数不同，可参数化，适合做成Skill。例如频繁回复的标准咨询，对应FAQ自动回复Skill。

\begin{funbox}
\textbf{居里夫人实验笔记中的蒸馏智慧。} 

玛丽·居里在提炼镭的过程中，从数吨沥青铀矿中反复溶解、沉淀、结晶、分离，历经四年才得到0.1克纯氯化镭。这个过程与Skill蒸馏惊人地相似：居里夫人面对的是海量矿石(大量AI对话历史)，通过反复的观察、归纳(提取成功模式)，最终凝练出纯净的镭元素(可复用的Skill模板)。她曾说：“生活中没有什么可怕的事物，只有需要理解的事物。”这句话同样适用于Skill蒸馏：不要把成功的AI交互看作一次性的幸运，而要主动去理解它为什么成功，从中提取出可复用的规律。
\end{funbox}

\subsection{Skill经典案例与设计实践}

针对不同角色和场景需求，小林发现Skill的设计思路各有侧重，下面从个人、团队和教学等角度分别展开。
\subsubsection*{\textbf{•个人场景：实验周报与文献精读}}
以小林为例，他设计了文献精读Skill：输入论文PDF，Skill自动提取摘要、方法、关键数据和结论，输出结构化的文献笔记。使用这个Skill后，小林阅读一篇论文的时间从两小时缩短到了40分钟。

小林的室友小杨，他每周都需要整理实验进展。他为自己设计了一个实验周报生成Skill，输入本周的实验完成情况和关键数据，自动生成结构化的周报。使用这个Skill之后，小杨的周报撰写时间从每周45分钟缩短到了5分钟。


\subsubsection*{\textbf{•团队场景：API文档与试题编制}}
对于研发团队而言，API文档自动生成Skill是另一类实用案例。输入源代码注释和函数签名，Skill自动生成格式统一的API参考文档，包括函数说明、参数列表、返回值类型和代码示例。



对于教学团队而言，试题编制Skill能根据教学大纲自动生成练习题。输入知识点列表和难度分布，Skill输出完整的试题及答案解析。




\subsubsection*{\textbf{•三类Skill卡片}}

Skill按功能可分为写作类、分析类和操作类三大类别，每类面向不同的场景需求，如图\ref{fig:ch3_3_16}所示。

写作类Skill包括邮件润色Skill和社交媒体文案Skill。邮件润色Skill输入一封草稿，输出三种风格版本(正式版、友好版、紧急版)，用户只需要写核心内容，Skill自动适配不同沟通场景。社交媒体文案Skill输入产品信息和目标受众，输出适配不同平台的文案风格(小红书活泼风、公众号专业风、微博短平快风)。



分析类Skill包括A/B测试分析Skill和用户反馈总结Skill。A/B测试分析Skill输入两组实验数据，输出统计显著性检验结果和业务建议，AI自动计算p值、置信区间，并给出可操作的解读。用户反馈总结Skill输入用户评论或访谈原始记录，输出关键词提取、情感分析和行动项。



操作类Skill包括数据清洗Skill和会议纪要Skill。数据清洗Skill输入原始表格(可能包含缺失值、格式不统一、重复记录)，输出经过清洗的标准化数据。会议纪要Skill输入会议录音文字记录，输出结构化的会议纪要(议题、讨论要点、结论、待办事项及责任人)。




\begin{figure}[H]
\centering
\BookFigureImage{./images_chapter_01-03/3.16.png}
\caption{三类Skill卡片}
\label{fig:ch3_3_16}
\end{figure}


\subsubsection*{\textbf{•组合Skill与迭代优化}}
\textbf{\textbf{组合Skill}\footnote{组合Skill：将多个子Skill串联起来处理复杂任务的高级用法。}：更高级的用法，}
将多个Skill串联形成组合Skill来处理复杂任务。比如市场调研报告Skill由三个子Skill组成。多个Skill串联形成组合Skill来处理复杂任务。以市场调研报告Skill为例，它由三个子Skill组成。第一个是行业数据搜索Skill，搜索最新的行业报告、市场数据和新闻。第二个是数据可视化Skill，将搜索到的数据转化为图表。第三个是报告撰写Skill，将分析结果和图表整合为结构化报告。

这三个Skill以流水线方式工作：搜索$\rightarrow$分析$\rightarrow$写作，前一个的输出自动成为后一个的输入。 如图\ref{fig:ch3_3_17}所示，多个子Skill可以串联为流水线，前一个Skill的输出自动成为后一个的输入，协作完成复杂任务。

\begin{figure}[H]
\centering
\BookFigureImage{./images_chapter_01-03/3.17.png}
\caption{组合Skill流水线}
\label{fig:ch3_3_17}
\end{figure}

优秀的Skill不是一次写好的，而是在实际使用中不断迭代完善的。
 小杨最初设计的实验周报Skill经过了四次迭代。v1.0只要求时间、实验内容、结果三栏；v2.0增加关键数据指标模块(用户反馈缺少定量数据)；v3.0增加结果解读和下周实验计划(用户反馈需要前瞻性分析)；v4.0绑定实验数据表实现自动数据抓取(用户反馈手动输入太麻烦)。
每次迭代都基于真实使用中的反馈，例如上次生成的周报少了数据对比，于是改进提示词；这次好了，但格式不够统一，继续改进。几个周期后，Skill的质量会有质的飞跃。

\subsection{Skill封装与调用机制}

打磨好了一个Skill，下一步是如何包装它，让别人也能方便地调用。

\subsubsection*{\textbf{•封装的三个要素}}

封装的三个要素如下。一是指令封装，将提示词固化，用户不需要了解内部的提示词细节，只需要输入核心参数，这让非技术用户也能轻松使用高级Skill。二是工具绑定，将Skill需要的外部工具与Skill关联，调用时自动启用，用户不需要手动切换工具。三是接口定义，明确定义输入(需要什么参数、什么格式)和输出(返回什么结果、什么格式)，一份好的接口定义是Skill可用性的关键。



\begin{figure}[H]
\centering
\BookFigureImage{./images_chapter_01-03/3.18.png}
\caption{Coze平台Skill配置界面}
\label{fig:ch3_3_18}
\end{figure}


\subsubsection*{\textbf{•不同平台的封装方式}}

在\textbf{Coze}中，创建一个Skill需要定义：触发词(用户说什么会触发这个Skill)、输入参数(参数名、类型、是否必填、默认值、示例值)、输出格式(文本/JSON/表格)、关联工具(搜索/计算/API)。 如图\ref{fig:ch3_3_18}所示，Coze平台的Skill配置界面包含触发词、输入参数、输出格式和关联工具等关键设置项。


在\textbf{Dify}中，Skill封装为工具类型：配置提示词模板(用\{\{变量名\}\}标出占位符)、关联知识库(可选)、设置输出变量、配置错误处理逻辑。 如图\ref{fig:ch3_3_19}所示，Dify平台的工具配置界面支持提示词模板编辑、知识库关联和输出变量设置。

\begin{figure}[H]
\centering
\BookFigureImage{./images_chapter_01-03/3.19.png}
\caption{Dify平台工具配置界面}
\label{fig:ch3_3_19}
\end{figure}


\subsubsection*{\textbf{•参数化调用的价值}}
\textbf{参数化\footnote{参数化调用：通过配置不同参数，让同一个Skill适配不同场景和需求的调用方式。}的价值 } 好的Skill设计支持参数化调用。以会议纪要Skill为例，它包含以下可配置参数：

会议类型(日报/周会/项目评审/季度回顾)；是否需要行动项列表(是/否)；是否需要待办事项责任人(是/否)；输出格式(Markdown/JSON/纯文本)；详细程度(简洁/标准/详细)。
两个用户调用同一个Skill，输入不同的参数值，得到不同风格的输出：一份封装，多种输出。这就是参数化的价值。
\begin{funbox}
\textbf{从角色扮演游戏技能树到AI技能模块。} 

Skill这个概念在AI领域的借用，很大程度上得益于角色扮演游戏(RPG)的普及。1985年的游戏《勇者斗恶龙》首创了角色技能系统，1990年代的《暗黑破坏神》将其发展为庞大的技能树：玩家通过解锁一个个技能节点，获得新的能力。这种将能力模块化、可组合的思想，与今天AI领域的Skill封装有着惊人的相似。只不过在游戏中，你解锁的是火球术和治愈术；在AI中，你解锁的是文献综述Skill和数据清洗Skill。就连组合Skill(多个子Skill串联)的使用方式，也和游戏中的连招如出一辙。
\end{funbox}

\subsection{Skill库与管理策略}

当小林拥有了十几个、几十个Skill之后，如何管理和查找就成了新问题。没有一个好的管理体系，Skill越多反而越混乱。


\subsubsection*{\textbf{•Skill的分类体系}}

Skill可以根据不同维度进行分类。按领域可分为写作类、分析类、操作类和决策类。按复杂度可分为基础Skill(单步骤)、复合Skill(多步骤)和工作流(含分支和条件判断)。按使用频率可分为高频(每天都用)、中频(每周用几次)和低频(偶尔使用)。在实际的Skill库管理中，通常会结合多个分类维度来组织，例如，一个周报生成Skill可同时归入“写作类”和“高频使用”两个分类标签下，便于用户从不同角度快速检索和调用。 如图\ref{fig:ch3_3_20}所示，Skill库支持按领域、复杂度和使用频率等多维度分类管理和卡片式展示。

\begin{figure}[H]
\centering
\BookFigureImage{./images_chapter_01-03/3.20.png}
\caption{Skill库管理界面}
\label{fig:ch3_3_20}
\end{figure}


\subsubsection*{\textbf{•版本管理与回退机制}}
Skill也需要像软件一样管理版本。v1.0$\rightarrow$v2.0$\rightarrow$v3.0，每次变更记录更新日志：新增了什么参数、优化了哪些提示词、修复了什么bug。当某个新版本效果不佳时，可以快速回退到旧版本。例如，一个竞品分析Skill最初版本只支持文本表格输出，用户反馈希望增加自动图表功能后，开发者在v2.0中集成了数据可视化模块，在v3.0中进一步加入了多语言支持。每个版本的更新日志都会详细记录变更内容、测试结果和回退方案，确保团队协作的顺畅。

\subsubsection*{\textbf{•分享协作与生态建设}}
在团队中共享Skill可以统一工作标准。
更广泛地，Coze的插件市场允许用户分享和安装他人创建的Skill，形成了一个Skill生态。这种生态极大地降低了入门门槛：不需要从零开始，先使用别人做好的Skill，再根据自己的需求进行修改和定制。

以下是蒸馏到调用的完整流程：
第一步：回顾与归纳。回顾过去10次文献检索对话，发现每次都要求：搜索PubMed、限定近一周、筛选与肿瘤免疫相关的文章、提取每篇的PICO信息、用表格呈现。这些共性就是Skill的雏形。
第二步：封装为Skill。
在Dify或Coze平台中将提示词、工具和流程封装在一起。提示词规定了搜索和分析的框架，工具连接学术搜索引擎，流程定义了从搜索到输出的完整链路。
第三步：调用Skill。以下展示调用Skill后的AI输出效果。

\begin{userbox}
调用每周文献简报Skill。
主题：肿瘤免疫治疗
关键词：CAR-T、免疫检查点抑制剂
筛选标准：影响因子大于5的期刊
\end{userbox}

\begin{aibox}
\textbf{本周肿瘤免疫方向新增文献简报}

\begin{formalTable}[H]
\caption{本周新增文献简报}
\label{tab:ch3_05_literature}
\begin{tabular}{c c c c}
\toprule
标题 & 期刊 & 核心发现 & 推荐 \\ \midrule
新型CAR-T设计 & Nature Med & 双特异性CAR-T提升实体瘤效果 & 五星 \\
PD-L1新机制 & Cell & 发现PD-L1的核内新功能 & 四星 \\
联合治疗策略 & JCO & 化疗+免疫联合方案生存获益 & 四星 \\ \bottomrule
\end{tabular}
\end{formalTable}

推荐阅读：本周重点推荐第1篇和第3篇。
\end{aibox}

这就是Skill的力量：不需要每周重复设计提示词，它已经帮助完成了从搜索到筛选到呈现的全流程。

\begin{center}
\textcolor{gray!70}{\*\ \ \*\ \ \*}
\end{center}
\vspace{12pt}

% ============================================================


\subsubsection*{小结}

Skill是封装了提示词、外部工具和执行流程的可复用能力单元。通过Skill蒸馏三部曲（回顾→归纳→泛化），可以将日常成功的AI交互经验固化为通用模板。多个子Skill还可以串联为流水线，前一个Skill的输出自动成为后一个的输入，协同完成过去需要数小时的复杂任务。

\bigskip
\subsection*{课后习题}

\begin{enumerate}
\setlength{\itemsep}{6pt}
\item 什么是Skill？它由哪三大要素组成？
\item 简述Skill蒸馏的三步法，并说明每一步的核心任务。
\item 组合Skill是如何工作的？请举一个将多个子Skill串联完成复杂任务的例子。
\item 参数化调用的价值是什么？同一个Skill配置不同参数如何适配不同场景？
\item Coze和Dify等平台分别如何封装和调用Skill？
\end{enumerate}

\newpage

\section[智能体规划执行闭环]{\textbf{智能体}\footnote{智能体(Agent)：能自主感知环境、制定计划、调用工具并执行任务的AI系统。}规划执行闭环}
% ============================================================

\begin{sectionkeypoints}
\item 理解智能体从被动响应升级为主动执行的核心转变，掌握感知、规划、执行、反馈的完整闭环
\item 掌握智能体感知阶段的三层解析过程(字面解析、意图推断、上下文融合)和主动追问策略
\item 理解任务分解的核心方法(深度优先、广度优先、动态规划)，掌握粒度平衡与失败处理策略
\item 熟悉工具调用的完整决策过程，掌握搜索、计算、API、文件、数据库等工具类型的适用场景
\item 了解多智能体系统的三种组织模式(主从、议会、市场)及冲突处理机制
\end{sectionkeypoints}




期刊投稿季，小林有一堆任务要处理：确定论文选题、检索目标期刊的投稿要求、整理实验数据、撰写Introduction部分。优秀的同学不会等到截稿日前一周才手忙脚乱。他们会先搞清楚目标期刊的偏好方向，然后制定写作计划，优先攻克最难的部分，完成后自我检测，根据审稿标准调整。
智能体的工作方式与此类似：给定一个目标，它自主完成从理解到执行的全过程。它不只是回答问题，而是完成任务。
\subsection{智能体感知与输入阶段}

智能体工作的第一步，是理解用户真正想要什么，小杨在尝试智能体时首先体会到这一点。这不是简单的听指令，而是一个主动解析的过程。 

感知的本质是从自然语言指令中提取结构化信息：目标是什么、约束条件有哪些、优先级如何、可用资源是什么、成功标准是什么。
如图\ref{fig:ch3_3_21}所示，智能体感知阶段将自然语言指令逐步解析为结构化信息，包括目标、约束、优先级和资源等要素。
\begin{figure}[H]
\centering
\BookFigureImage{./images_chapter_01-03/3.21.png}
\caption{智能体感知阶段流程}
\label{fig:ch3_3_21}
\end{figure}



\vfill
\begin{knowledgebox}

\subsubsection*{\textbf{•感知的三层解析}}
\textbf{从愿景到现实：智能体如何理解用户意图?} 

1995年，MIT的帕蒂·梅斯在《麻省理工科技评论》上描绘了软件智能体的愿景：它们是能感知环境、做出决策、代表用户执行任务的自主程序。三十年后，这个愿景进化为我们今天看到的AI Agent。但要让智能体真正理解用户，需要经过三个层次的解析：
第一层：字面解析，即提取关键词和句法结构。“帮我找一家便宜的酒店”$\rightarrow$\{动作：找，对象：酒店，属性：便宜\}
第二层：意图推断，即基于常识补全缺失信息。“便宜的酒店”需要知道目的地、日期才能执行$\rightarrow$识别出信息缺口
第三层：上下文融合，即结合对话历史、用户画像、环境信息修正理解。如果用户刚说过“下周去北京出差”，智能体自动将目的地推断为“北京”
这三层依次递进：只有完成每一层的解析，智能体才能从“听懂你说了什么”升级到“理解你想要什么”。
\end{knowledgebox}
\vspace{1em}

感知的具体案例如下。当用户说“帮我找一家便宜的酒店”，智能体需要解析出目标(预订酒店)、约束(价格低)，同时识别出缺失的信息即目的地和日期，需要追问用户。第二个案例，用户说“分析一下这份销售数据”，智能体需要判断数据类型、用户可能关心的维度、期望的分析深度。


\subsubsection*{\textbf{•主动追问策略}}
那么，主动感知在实践中是如何工作的呢?当用户说“帮我找一家便宜的酒店”时，优秀的智能体不会直接去搜索“便宜的酒店”，因为这个指令太模糊了。它会主动追问：“好的，我来帮您找酒店。请问您想去哪个城市?计划什么时间入住?预算大概是多少?您对酒店有什么特殊要求吗(如靠近地铁、含早餐、有无烟房)?”


这种追问不是随机的，而是基于对任务结构的理解，因为订酒店需要知道目的地、时间、预算、偏好这四个核心信息。智能体识别到缺失的信息后，逐一询问。等用户回答完毕，它才开始执行。

主动感知与被动感知的区别在于信息获取方式。被动感知要求用户一次性给出完整、明确的指令，这对用户的提示词编写能力要求较高。主动感知则是智能体在信息不足时主动追问，这是优秀智能体的重要特征，也是它与普通问答系统的核心区别。在主动感知模式下，智能体会检测到关键信息缺失，主动向用户询问以补充完善。
如图\ref{fig:ch3_3_22}所示，当用户指令信息不足时，智能体会主动追问以补充缺失信息，确保任务执行的准确性。

\begin{figure}[H]
\centering
\BookFigureImage{./images_chapter_01-03/3.22.png}
\caption{智能体主动追问示例}
\label{fig:ch3_3_22}
\end{figure}


\subsubsection*{\textbf{•多模态感知能力}}

多模态感知方面，智能体接收的不只是文字，它可以看到图片、听到语音、读取文档。例如一个穿搭顾问Agent接收用户上传的衣柜照片，识别衣物种类和颜色，根据天气和场合给出穿搭建议。一个数据分析Agent读取用户上传的Excel文件，识别列名和数据类型，自动选择合适的分析方法。一个合同审查Agent接收PDF合同文件，提取关键条款，与标准条款对比，标注差异和风险点。 

\begin{funbox}
\textbf{搜索引擎如何读懂你的心。} 

1998年，PageRank算法革命性地改变了网页排序方式：它不是看网页内容说了什么，而是看其他网页如何链接它。这种通过外部链接判断内容价值的思想，后来被应用到智能体感知中：当用户给出模糊指令时，智能体不是只看字面意思，而是通过上下文(对话历史、用户画像、环境信息)来推测真实意图。就像早期的搜索引擎遇到苹果这个词不知道是水果还是公司，智能体也需要多维度信息才能做出准确的意图判断。今天的大模型已经能根据对话历史自动补充缺失信息，但主动追问仍然是保证执行准确性的黄金法则：永远不要假设AI已经猜到了你的意思。
\end{funbox}




\subsection{任务规划与过程拆解}

理解了目标之后，智能体需要回答一个关键问题，小杨的项目经验让他深刻理解了这个环节：先做什么，再做什么?
这就是任务分解：将大目标拆解为有依赖关系的子任务。
掌握了目标拆解的思路后，智能体需要选择执行路径。

\begin{figure}[H]
\centering
\BookFigureImage{./images_chapter_01-03/3.23.png}
\caption{智能体任务分解树形图}
\label{fig:ch3_3_23}
\end{figure}


\subsubsection*{\textbf{•新能源汽车行业报告完整案例}}
以“帮我写一份新能源汽车行业报告”为例，智能体的规划可能是：

\[
\text{目标} \rightarrow \{\text{任务}_1， \text{任务}_2， \dots， \text{任务}_n\}
\]

其中，
任务\(_1\)：搜索最新的行业数据和市场规模(搜索工具)；任务\(_2\)：分析头部企业的竞争格局(基于任务\(_1\)的结果)；任务\(_3\)：梳理政策环境变化和影响(搜索工具)；任务\(_4\)：生成报告框架(基于所有前置任务)；任务\(_5\)：逐节撰写内容(基于框架和各子任务结果)；任务\(_6\)：格式美化和引用检查(最终质检)。

假设智能体在执行过程中遇到了如下情况：搜索后发现2025年行业市场规模数据在公开渠道没有完整版本。智能体自动将任务1拆分为两个子任务，任务1a：搜索已公开发布的分项数据(新能源乘用车销量、充电桩数量等)；任务1b：基于分项数据用统计模型估算整体市场规模，并标注估算值及置信度。任务2则基于1a和1b的结果展开竞争格局分析，在最终报告中标注数据来源的可靠性等级。如图\ref{fig:ch3_3_23}所示，展示了智能体任务分解的过程。


\subsubsection*{\textbf{•三种规划策略}}

深度优先是按顺序依次完成每个子任务，做完一个再做下一个，适合子任务之间强依赖的场景。广度优先是同时推进多个没有依赖关系的子任务，适合需要并行收集信息的场景。动态规划是根据执行结果实时调整后续子任务的方案，适合不确定性高的复杂任务。 如图\ref{fig:ch3_3_24}所示，深度优先策略按顺序逐一完成子任务，广度优先策略则同时推进无依赖关系的子任务，两者适用于不同场景。

\begin{figure}[H]
\centering
\BookFigureImage{./images_chapter_01-03/3.24.png}
\caption{深度优先与广度优先策略对比}
\label{fig:ch3_3_24}
\end{figure}


\subsubsection*{\textbf{•粒度平衡与失败处理}}
规划中的两个关键挑战：
一是粒度问题，拆得太粗则子任务依然复杂难执行，拆得太细则规划本身耗费大量计算资源，好的规划应该在粒度适中和可执行性之间找到平衡，一般以一个子任务可以在1-2步内完成为参考标准。二是失败处理，当某个子任务无法完成时(搜索API无返回、数据源不可用)，智能体需要能随机应变，方案一失败后自动切换到方案二，而不是卡住不动或报错。



\subsection{工具调用与执行环节}

小杨在使用智能体时发现，LLM本身只有语言能力，它不能搜索实时信息、不能执行精确计算、不能操作文件。要让智能体真正做到做事，它需要调用外部工具。 


\subsubsection*{\textbf{•常见工具类型与场景}}

网络搜索适用于需要实时信息的场景，比如最新新闻、当前股价、天气情况。没有搜索工具，AI的知识截止于训练数据的日期。

计算器或代码执行器适用于需要精确计算的场景，比如统计分析、数据计算、公式推导。LLM的数学能力有限，复杂计算必须借助工具。

API连接器适用于调用第三方服务的场景，比如查航班、订酒店、发邮件、发短信。通过API，AI能操作现实世界中的系统。

文件读写适用于处理文档的场景，比如读取PDF、Word或Excel，生成报告、保存数据。没有文件工具，AI只能说不能写。

数据库查询适用于检索和更新数据的场景，比如查客户信息、更新订单状态、生成数据报表。如图\ref{fig:ch3_3_25}所示，绑定搜索、计算、API等外部工具后，AI的能力边界从纯语言处理扩展到实时信息获取和精确计算。

\begin{figure}[H]
\centering
\BookFigureImage{./images_chapter_01-03/3.25.png}
\caption{绑定工具前后的AI能力对比}
\label{fig:ch3_3_25}
\end{figure}


\subsubsection*{\textbf{•工具调用的决策链路}}
智能体如何知道什么时候用什么工具?
以“查一下明天北京到上海的机票”为例，
(1) 智能体判断：这是一个实时信息查询，需要调用航班查询API，(2) 解析参数：出发地=北京、目的地=上海、日期=明天，(3) 调用航班API，传入解析后的参数，(4) 解析API返回的JSON数据(航班号、时间、价格、余票等)，(5) 将原始数据转化为用户友好的自然语言回答。 如图\ref{fig:ch3_3_26}所示，智能体工具调用经历指令解析、参数提取、API调用和结果整合的完整链路。

\begin{figure}[H]
\centering
\BookFigureImage{./images_chapter_01-03/3.26.png}
\caption{工具调用全链路}
\label{fig:ch3_3_26}
\end{figure}


工具调用不是每次都能成功，优秀的智能体需要处理异常。例如搜索无结果时更换搜索词或切换搜索引擎；API超时时重试，重试失败则提示用户手动操作；数据格式异常时尝试其他解析方式或请求用户确认。除了自我检查，智能体还需要处理来自外部的反馈信号，包括工具返回的错误信息、用户的中途干预、环境状态变化等。

\subsection{反馈调整与动态优化}



\subsubsection*{\textbf{•自我反思机制}}
自我反思是指智能体主动检查自己的输出，发现问题并修正。

智能体的第一次执行结果往往不完美。 小张让智能体帮他整理上周的实验数据报告。第一次执行的结果中，数据表格的格式混乱、图表的坐标轴标注不清。具体来说：表格里的p值写成了0.05而没有标明具体数值，箱线图的纵轴标注为值而非具体的抗拉强度(MPa)。智能体自我检查后，在第二轮输出中：将p值改写为F=28.6， p<0.001的标准统计格式，将箱线图纵轴改为抗拉强度(MPa)，并增加误差棒。第三次检查发现还缺少对照组样本量的说明，又补充了n=5 per group的注释。这个自我修正过程仅用了三次迭代。关键不在于第一次就做对，而在于快速发现错误并自我修正。




\subsubsection*{\textbf{•外部反馈信号}}
除了自我检查，智能体还需要处理来自外部的反馈信号：工具返回的错误信息、用户的中途干预、环境状态的变化。 

闭环优化的典型场景：
AI编程助手写完代码后执行测试，测试失败则分析错误日志并修改代码，然后重新测试，直到通过。AI内容创作写完文章后由另一个AI做质量评估，如果评分低于阈值则修改并重新评估，直到达标。AI数据分析生成图表后检查数据准确性，发现数据映射错误则修正参数并重新生成。如图\ref{fig:ch3_3_27}所示，反馈闭环通过执行、检查、修正、再执行的迭代过程持续优化输出质量。

\begin{figure}[H]
\centering
\BookFigureImage{./images_chapter_01-03/3.27.png}
\caption{反馈闭环循环}
\label{fig:ch3_3_27}
\end{figure}


\subsubsection*{\textbf{•反馈频率的平衡艺术}}
过于频繁的检查会降低效率；过于稀少的检查可能让错误累积到无法修复的程度。好的策略是在关键节点设置检查点，高风险操作前强制检查。

 以下是智能体闭环的完整演示：

\textbf{第一步：感知阶段}

\begin{userbox}
帮我分析这组合金材料的力学性能测试数据。
数据文件已上传包含三组样品：合金A、合金B、合金C。
每组测试了抗拉强度、屈服强度、延伸率和硬度四个指标。
每组有5个平行样。
\end{userbox}

\begin{aibox}
收到!我来帮您分析。在开始之前我需要确认：
1. 是否需要做统计分析(如均值、标准差、显著性检验)?
2. 是否需要生成可视化图表(如柱状图、箱线图)?
3. 是否需要与行业标准或文献数据进行对比?
4. 报告的输出格式偏好(Word/PDF/LaTeX)?
\end{aibox}

可以看到智能体首先主动追问来补充缺失的信息(分析深度、输出格式等)，而不是拿着不完整的指令盲目开始。 小杨在描述需求时只说了分析数据，但智能体识别到分析深度图表类型数据对比基准输出格式这四个关键信息缺失，因此逐一追问补齐。如果小杨此时补充说需要做显著性检验，智能体还会进一步追问显著性水平设为多少，这种递进式追问确保了最终的分析结果完全符合用户的真实需求。

\textbf{第二步：规划阶段}

智能体内部将任务拆解为：
(1) 读取和验证上传的数据文件，(2) 计算各指标的描述性统计量(均值、标准差、变异系数)，(3) 进行组间差异的统计检验(方差分析)，(4) 生成可视化图表，(5) 撰写分析报告。

\textbf{第三步：工具调用与执行}

\begin{aibox}
\textbf{数据分析结果如下：}

\textbf{描述性统计}

\begin{formalTable}[H]
\caption{合金样品力学性能描述性统计}
\label{tab:ch3_06_materials}
\begin{tabular}{c c c c c}
\toprule
样品 & 抗拉强度 & 屈服强度 & 延伸率(\%) & 硬度 \\ \midrule
合金A & 856$\pm$32 & 721$\pm$28 & 12.3$\pm$1.2 & 42$\pm$2 \\
合金B & 923$\pm$28 & 784$\pm$25 & 8.5$\pm$0.9 & 48$\pm$3 \\
合金C & 798$\pm$35 & 665$\pm$30 & 15.1$\pm$1.5 & 38$\pm$2 \\ \bottomrule
\end{tabular}
\end{formalTable}

\textbf{方差分析结果}

抗拉强度：F=28.6 p$<$0.001(组间差异极显著)
延伸率：F=42.3 p$<$0.001(组间差异极显著)
硬度：F=35.1 p$<$0.001(组间差异极显著)

\textbf{结论}

合金B强度最高但延伸率最低。合金C强度最低但塑性最好。合金A综合性能最优。实际应用中应根据具体需求选择。
\end{aibox}

\textbf{第四步：反馈调整 }

 在第一步输出分析结果后，智能体自动进行了质量检查，发现报告中缺少了测量不确定度分析这一关键环节，它会自我检查发现并修正：

\begin{aibox}
注意：报告中缺少了测量不确定度分析。正在补充计算。
补充结果：各测量值的不确定度均在$\pm$2\%以内，符合ASTM E8标准要求。
\end{aibox}

\subsection{多智能体与协作分工}

为什么需要多智能体?复杂任务往往需要多个不同专业能力的角色协作。就像软件开发需要产品经理、架构师、工程师、测试员各司其职，一个智能体很难同时胜任所有角色。

当一个任务太复杂，一个智能体搞不定时怎么办?答案是：用多个智能体，分工协作。 如图\ref{fig:ch3_3_28}所示，多智能体系统采用主从、议会和市场三种组织模式，分别适用于不同复杂度的协作场景。

\begin{figure}[H]
\centering
\BookFigureImage{./images_chapter_01-03/3.28.png}
\caption{多智能体协作架构}
\label{fig:ch3_3_28}
\end{figure}

\subsubsection*{\textbf{•三种组织模式}}





主从模式是一个主Agent(协调者)调度多个子Agent(执行者)，主Agent负责任务分解、分发结果、整合输出。议会模式是多个Agent讨论、辩论、达成共识，适合需要综合多方意见的决策任务。市场模式是Agent自主发布任务和认领，适合开放式、不确定性高的任务分派。


\subsubsection*{\textbf{•多智能体通信方式}}
Agent之间通过结构化消息交换中间结果。例如多智能体系统中，CEO Agent确定需求，产品经理Agent输出PRD文档，设计师Agent产出UI原型，程序员Agent编写代码，测试员Agent检测Bug并反馈给程序员修改。在多无人机协作中，每个无人机是一个Agent，通过通信网络共享位置信息和任务状态，协作完成搜索和救援任务。

当多个Agent对同一问题给出不同结论时，可以引入协调Agent综合双方论据，或引入第三方验证Agent重新核查。

以上讨论的智能体都生活在数字世界中。但如果智能体要操作物理世界，例如控制机械臂、驾驶汽车、操作手术刀，就会面临全新的挑战：实时性、安全性、物理约束。同时，无论虚拟还是物理，智能体的权限边界在哪里?谁该为此负责?这引出了本章的最后一节。



\subsubsection*{\textbf{•期末考试助手多智能体系统搭建案例}}

纸上得来终觉浅，本节以Dify平台为例，带领读者动手搭建一个由三个智能体组成的“期末考试论文写作助手”。通过这个完整的操作案例，你会直观感受到三种组织模式（主从模式）在实际平台中的落地方式，以及智能体之间如何分工协作。

\textbf{准备工作：}在浏览器中打开Dify官网，注册账号后进入工作区。Dify提供了可视化的工作流编排界面，无需编写代码即可完成多智能体系统的搭建。

\textbf{第一步：创建三个智能体}

在Dify的“工作室”页面，依次创建以下三个智能体：

\begin{enumerate}
\item \textbf{文献搜索Agent}（主模式）：系统指令设为“你是一名学术文献检索助手，负责根据给定主题搜索最新、最相关的学术文献。请优先调用搜索工具获取实时数据，并按APA格式输出文献信息。”在工具设置中为其绑定“网络搜索”工具。

\item \textbf{数据分析Agent}（主模式）：系统指令设为“你是一名数据分析师，负责对提供的实验数据或文献数据进行统计分析，输出均值、标准差、显著性检验结果，并生成可视化图表。”在工具设置中为其绑定“代码执行器”工具。

\item \textbf{论文写作Agent}（主模式）：系统指令设为“你是一名学术论文写作助手，负责根据文献搜索和数据分析的结果，撰写结构清晰、语言规范的课程论文。注意保持学术性，区分自己的观点和引用的观点。”该Agent不需要绑定工具，专注文本生成。
\end{enumerate}

\textbf{第二步：编排工作流（搭建主从模式）}

创建完三个智能体后，进入Dify的“工作流”编排界面，通过拖拽方式将它们连接起来，形成流水线式协作。具体操作如下：

\begin{enumerate}
\item 从左侧节点面板拖入一个“开始”节点，设定输入变量为\texttt{topic}（论文主题），数据类型为“文本”。

\item 拖入一个“LLM节点”，命名为“文献搜索”，选择刚创建的文献搜索Agent。在提示词中引用输入变量：\texttt{请搜索关于\{\{topic\}\}的最新学术文献，重点关注近五年（2020-2025年）的高被引论文。}

\item 拖入第二个“LLM节点”，命名为“数据分析”，选择数据分析Agent。其输入引用上一步的输出：\texttt{基于以下文献信息，分析该领域的研究趋势：\{\{文献搜索Agent的输出\}\}}。

\item 拖入第三个“LLM节点”，命名为“论文写作”，选择论文写作Agent。其输入同时引用前两步的结果：\texttt{请根据以下文献列表和数据分析结果，撰写一篇关于\{\{topic\}\}的课程论文初稿。文献列表：\{\{文献搜索Agent的输出\}\}。数据分析：\{\{数据分析Agent的输出\}\}}。

\item 最后拖入“结束”节点，将论文写作Agent的输出作为最终结果返回。
\end{enumerate}

\textbf{第三步：运行测试}

保存工作流后，点击“运行”按钮，输入论文主题，例如“深度学习在肺结节CT诊断中的应用”。系统将依次执行三个智能体：

\begin{aibox}
\textbf{运行日志：}

文献搜索Agent正在执行... 完成。搜索到12篇相关文献，其中高被引论文5篇。 关键文献：Liu et al. (2024) Medical Image Analysis、Wang et al. (2023) IEEE TMI...

数据分析Agent正在执行... 完成。趋势分析：该领域年发文量从2020年的120篇增长至2024年的480篇。 常用模型：CNN占比42\%，U-Net占28\%，Transformer占20\%。

论文写作Agent正在执行... 完成。已生成含摘要、引言、方法、结果、结论的完整论文初稿，约3500字。
\end{aibox}

从运行日志可以清晰看到三个智能体的协作流程：文献搜索Agent先完成检索，将其结果传递给数据分析Agent，数据分析Agent完成趋势分析后再传递给论文写作Agent，最终产出一篇结构完整的课程论文初稿。整个过程完全自动执行，用户仅需输入论文主题。

\textbf{多智能体协作的关键观察：}

通过这个动手案例，你可以直观地观察到以下几点：

第一，这就是主从模式的典型实现：用户（主控制器）设定目标后，由工作流引擎依次调度三个子Agent，各自完成分配的任务，最终由主控制器整合输出。

第二，智能体之间的通信是通过结构化数据传递的：前一节点的输出自动转化为后一节点的输入变量，参数值的错误或格式不匹配会直接导致下游任务失败，因此需要像设置接口一样仔细定义每个Agent的输入输出规范。

第三，如果某个Agent执行失败（例如搜索工具返回超时），Dify的工作流引擎会自动触发重试机制，最多重试3次后仍失败则跳过该节点并标记为警告。这正是实践中“冲突处理”和“失败降级”的自动化实现。



通过这个案例，我们可以对多智能体系统的实际搭建有更深刻的切身体验，从三种组织模式的理解，到具体平台上的拖拽编排，再到冲突处理和失败降级的实际应对，层层递进地掌握整个搭建流程。




% ============================================================


\begin{center}
\textcolor{gray!70}{\*\ \ \*\ \ \*}
\end{center}
\vspace{12pt}



\subsubsection*{小结}

智能体通过感知、规划、执行和反馈调整四个环节构建自主执行闭环。感知阶段通过三层解析和主动追问补充缺失信息；规划阶段将大目标拆解为有依赖关系的子任务；执行阶段调用搜索、计算、API等外部工具完成具体操作；反馈阶段通过自我反思和外部信号持续优化输出质量。当单个智能体难以胜任复杂任务时，多智能体系统通过主从、议会或市场等组织模式实现分工协作。借助Dify等平台，用户只需通过拖拽即可搭建多智能体工作流，将概念落地为实际可运行的协作系统。

\bigskip
\subsection*{课后习题}

\begin{enumerate}
\setlength{\itemsep}{6pt}
\item 智能体的感知阶段需要经过哪三个层次的解析？各层次分别解决什么问题？
\item 任务规划中深度优先和广度优先两种策略各自适用于什么场景？
\item 智能体调用外部工具时，常见的工具类型有哪些？各适用于什么任务？
\item 反馈调整环节中，自我反思和外部反馈分别指什么？
\item 多智能体系统的三种组织模式（主从、议会、市场）各有什么特点？
\end{enumerate}

\newpage
\section{具身智能与安全边界}
% ============================================================

\begin{sectionkeypoints}
\item 了解具身智能的基本概念及其与虚拟AI的核心差异，理解实时性、鲁棒性和安全性三大独特挑战
\item 掌握物理闭环(感知$\rightarrow$规划$\rightarrow$控制)的核心机制，了解Sim-to-Real技术路径及仿真间隙挑战
\item 理解智能体权限分级模型和最小权限原则，掌握沙箱机制、人机确认和操作日志等安全管控手段
\item 了解数据隐私保护的核心技术(本地部署、数据脱敏、联邦学习、差分隐私)，理解具身智能的隐私风险特殊性
\item 掌握人机协作的五个层次(L0至L4)及其责任归属原则，理解可解释AI对建立信任的关键作用
\end{sectionkeypoints}




一家工厂引进了AI视觉检测系统来检查产品质量。一切顺利运行了三个月后，有一天系统突然将所有的合格品判定为次品，产线被迫停线。工程师排查后发现，原因很简单：车间更换了一组照明灯管，光线角度和色温变了，AI模型从未见过这种光照条件。
这个案例揭示了一个重要问题：当AI从屏幕走进物理世界，它会面临完全不同的挑战。同时，如果这个AI有权限删除数据、控制机械设备、操作银行账户，谁为它的错误负责?



\subsection{具身智能基本概念}

\textbf{具身智能}\footnote{具身智能：拥有物理实体、能在真实环境中感知和行动的AI系统，是AI从虚拟世界走向物理世界的核心技术方向。}是指拥有物理实体、能在真实环境中感知和行动的AI系统。简单说，就是有身体的AI。\footnote{Shakey机器人项目， SRI International， 1966-1972.}

\begin{funbox}
\textbf{第一个拥有身体的AI。} 

1966年，斯坦福国际研究所(SRI)启动了一个在当时看来近乎疯狂的项目：打造世界上第一个能够感知-推理-行动的移动机器人。这个被称为Shakey的机器人，装配了电视摄像头、三角测距仪和触觉传感器，通过无线电与一台巨大的计算机相连。Shakey能够规划路径、推开障碍物、按指令推动箱子。它在1972年被《Life》杂志评为年度最佳机器人，其核心软件系统STRIPS后来成为AI规划领域的经典算法。Shakey揭示了具身智能的一个基本真理：物理世界中的每一步行动，都需要感知、规划和控制三者的精密协作：这个真理至今仍然成立。
\end{funbox}


\subsubsection*{\textbf{•虚拟AI与具身智能的核心差异}}


虚拟AI的输入输出是信息，环境是数字化的，没有物理约束。发错一条消息可以撤回，算错一个数字可以重算。具身智能的输入输出涉及物理世界，环境是连续的、不确定的、有物理约束的(重力、摩擦力、惯性)。抓错一个零件可能导致产品报废，走错一步可能摔倒。这就是为什么具身智能系统的设计必须将物理安全保障作为最高优先级。一个虚拟AI算错了一个数学题，用户可以重算；但一个送货机器人算错了路线，它可能撞到行人或货架。


\subsubsection*{\textbf{•具身智能的三个能力支柱}}


具身智能的三个能力支柱包括感知、决策和控制。感知通过传感器获取环境信息：相机(视觉)、激光雷达(距离)、触觉传感器(接触)、力传感器(力度)。决策是理解环境状态并做出行动选择。控制是精确执行物理动作：移动、抓取、避障，每个动作需要在毫秒级精度内完成。

\subsubsection*{\textbf{•具身智能的三大独特挑战}}

具身智能面临的三大独特挑战是实时性、鲁棒性和安全性。实时性要求毫秒级响应，虚拟AI可以慢慢思考，但机器人不能，不及时的避障动作就是碰撞。鲁棒性要求传感器噪声、光照变化、摩擦力变化都不影响性能，必须在各种不确定条件下稳定工作。安全性要求错误动作不能导致物理伤害或财产损失，不能像虚拟AI那样错了就重来。在工业机器人领域，一次错误的焊接指令可能导致整条产线的产品报废，损失数十万元，这就是为什么工业机器人的AI系统必须经过数万小时的安全测试才能部署到产线上。


\subsubsection*{\textbf{•Sim-to-Real技术与仿真间隙挑战}}
\textbf{\textbf{Sim-to-Real}\footnote{Sim-to-Real(仿真到现实迁移)：在虚拟环境中训练的策略迁移到真实物理世界的方法。}(仿真到现实)：} 在真实物理世界中训练机器人成本高、风险大、速度慢，一次失败的操作可能导致设备损坏。因此研究人员先在虚拟仿真环境中训练，然后将学到的策略迁移到真实机器人上。 如图\ref{fig:ch3_3_30}所示，Sim-to-Real技术先在虚拟仿真环境中训练策略，再将其迁移到真实物理世界的机器人上。

\begin{figure}[H]
\centering
\BookFigureImage{./images_chapter_01-03/3.30.png}
\caption{Sim-to-Real技术路径}
\label{fig:ch3_3_30}
\end{figure}

Sim-to-Real的关键挑战：仿真间隙。虚拟环境永远无法完美模拟真实世界，因为光照不同、摩擦力不同、材料特性不同。解决这个问题的常用技术是域随机化，也就是在虚拟训练中随机改变光照条件、物体纹理、物理参数等，让模型在变化的环境中学习，从而提高它在真实世界中的适应能力。

\subsubsection*{\textbf{•里程碑案例}}


具身智能的里程碑案例包括以下几个。Figure 01(2024年)是Figure AI展示的通用人形机器人，在仓库中自主工作，当人类说“把那个箱子搬到对面的货架上”时，它通过视觉识别箱子位置和类型，自主规划搬运路径，调整抓取角度和力度。波士顿动力的Atlas能在复杂地形中保持平衡、跑酷、后空翻，2024年最新版本已能自主搬运重物。Google DeepMind的RT-2将视觉-语言模型直接映射到机器人控制指令，能理解和执行从未在训练数据中出现过的组合指令。

% 智能体在完成任务的过程中会接触到大量敏感信息，数据在哪里，风险就在哪里。智能客服Agent处理用户咨询时会接触到姓名、电话、地址、身份证号等个人敏感信息。医疗AI分析病历数据时必须确保数据不出院区。金融AI处理客户资产信息时需要符合严格的合规要求。

\subsection[\textbf{感知行动与物理闭环}]{感知行动与物理闭环\footnote{物理闭环(Physical Closed Loop)：具身智能系统中感知→规划→控制的持续循环。}}

机器人在物理世界中是如何工作的?核心机制是\textbf{物理闭环}：感知$\rightarrow$规划$\rightarrow$控制，每秒钟循环数十次到数百次。 如图\ref{fig:ch3_3_29}所示。

\begin{figure}[H]
\centering
\BookFigureImage{./images_chapter_01-03/3.29.png}
\caption{物理感知-规划-控制闭环}
\label{fig:ch3_3_29}
\end{figure}

\subsubsection*{\textbf{•物理闭环机制}}


物理闭环的核心在于感知、规划和控制三个环节的持续循环：传感器获取环境信息（感知），算法决定下一步动作（规划），执行器完成物理操作（控制），随后新的传感器数据再次反馈到系统。


\subsubsection*{\textbf{•无人机飞行与机械臂抓取案例}}
无人机利用IMU(惯性测量单元)检测自身的倾斜角度，将这个角度输入控制算法，算法计算四个旋翼各自的转速调整量，电机执行调整，新的姿态被传感器检测到，再次反馈到算法。这个闭环每秒钟执行400次以上。


机械臂抓取鸡蛋的过程涉及多个传感器协同工作。首先，视觉系统要检测鸡蛋的位置、姿态和朝向(感知环节)。其次，规划系统要计算抓取轨迹(规划环节)。接着，控制系统驱动机械臂沿轨迹移动并调整夹爪开合度(控制环节)。同时，力传感器检测夹爪与蛋壳的接触压力(再次感知环节)。最后，根据压力反馈进行调整，压力太大则减小夹持力，太小则增加。

整个过程在几百毫秒内完成，涉及视觉、力觉、位置觉三种传感信息的实时融合。 这一过程的复杂性在于：如果压力检测延迟了50毫秒，夹爪可能已经挤压过度导致蛋壳破裂；如果视觉定位偏差了1毫米，夹爪可能错过目标。这就是具身智能与虚拟AI的本质区别：在数字世界，1毫米的偏差无足轻重；在物理世界，1毫米可能就是成败的分界线。




% \textbf{具身智能的三大独特挑战}





\subsection{智能体权限与安全管控}

当智能体能读写文件、发送邮件、调用API、控制设备时，如何防止它乱来?

\subsubsection*{\textbf{•权限分级模型}}

\textbf{只读级 }可读取信息，但不能修改。如查阅文档、搜索知识库、读取数据库；\textbf{执行级 }可执行预定义操作。如发送消息、查询数据、生成报告、触发工作流；\textbf{管理级 }可修改配置、安装工具、管理用户账号。此级别需要严格审核；\textbf{超级级 }可访问系统核心资源、修改系统设置。一般不直接赋予AI。 如图\ref{fig:ch3_3_31}所示，智能体权限分为只读级、执行级、管理级和超级级四个层级，越往上权限越大但管控也越严格。

\begin{figure}[H]
\centering
\BookFigureImage{./images_chapter_01-03/3.31.png}
\caption{智能体权限层级金字塔}
\label{fig:ch3_3_31}
\end{figure}


\subsubsection*{\textbf{•沙箱机制与最小权限原则}}
\textbf{沙箱机制：} \textbf{沙箱}\footnote{沙箱：为程序提供隔离执行环境的安全机制，防止对主系统造成影响。即使程序崩溃或执行恶意代码，也不会影响沙箱外的系统。}是为程序提供隔离执行环境的安全机制。智能体的沙箱机制可以限制其对文件系统、网络资源和系统调用的访问范围，确保即使AI执行了意外的操作，也不会影响系统核心功能和敏感数据。例如，一个文档处理Agent应当在沙箱环境中运行，只能访问指定的工作目录，无法读取系统配置文件或其他用户的私有数据。

\textbf{\textbf{最小权限原则}\footnote{最小权限原则(Principle of Least Privilege)：只给智能体执行任务所需的最小权限的安全策略。}：} 只给智能体完成任务所需的最小权限。一个文件整理Agent只需要读写指定工作目录下的文件，不需要访问系统目录。一个数据分析Agent只需要数据库的读取权限，不需要写入、修改或删除权限。一个邮件发送Agent只需要发送权限，不需要读取所有邮件的权限。


\subsubsection*{\textbf{•人机确认机制}}
\textbf{人机确认机制：高风险操作需要人类确认后才能执行}。例如“确认删除目录下的所有文件?【是/否】”；“已草拟了以下邮件，请确认后发送。【确认发送/取消】”；“检测到系统配置将被修改，当前操作需要管理员权限。”

% \textbf{隐私保护的核心技术：} 隐私保护的核心技术包括本地部署、数据脱敏、联邦学习和差分隐私。本地部署是将LLM部署在企业内部的服务器上，所有数据处理都在内网完成，数据不出防火墙。数据脱敏是在输入AI前自动识别并替换敏感信息。联邦学习使多个医院可以在不共享用户病历的前提下协同训练诊断模型。差分隐私是在模型输出或数据查询结果中注入受控的随机噪声。


\subsubsection*{\textbf{•操作日志与溯源审计}}
操作日志与审计是指记录AI的所有操作，生成不可篡改的审计日志。当出现问题时，通过日志回溯找到原因和责任人。

在某企业的AI运维系统中，审计日志记录了每个AI操作的完整链路：何时由哪个Agent调用了哪个API，传入了什么参数，返回了什么结果，耗时多久。当一次意外的数据库写入事件发生时，运维团队通过回溯审计日志在30分钟内就定位到了问题Agent和触发条件，而如果没有日志系统，排查类似问题通常需要数天时间。

\subsection{数据隐私与伦理风险}

上文讨论的权限管控针对的是智能体的操作能力边界，而数据隐私问题源于具身智能的另一独特性质：通过传感器持续采集物理世界的各类信息。虚拟AI的输入是用户主动提供的文本或文件，用户清楚自己给了什么数据；但具身智能通过摄像头、麦克风、激光雷达等传感器被动地获取环境信息。一台巡检机器人经过办公室时会录下员工交谈，一台家庭服务机器人会全天候拍摄室内画面。这些传感器数据包含了大量个人隐私信息：人脸、声音、行为模式，而数据被采集时用户甚至未必知晓。

智能体在完成任务的过程中会接触到大量敏感信息，如何保护隐私是一项无法回避的挑战。


\subsubsection*{\textbf{•合规要求}}

合规要求方面，中国《生成式人工智能服务管理暂行办法》(2023年8月15日施行)要求AI生成内容不得包含违法信息，服务提供者需对训练数据合法性负责。欧盟AI法案(2024年通过)\footnote{European Commission， The EU Artificial Intelligence Act， 2024.}将AI系统按风险分为四级：不可接受风险(禁止)、高风险(严格合规要求)、有限风险(透明度义务)、极低风险(无额外义务)。


\subsubsection*{\textbf{•隐私保护核心技术}}

\textbf{本地部署}：将LLM部署在企业内部的服务器上，所有数据处理都在内网完成，数据不出防火墙；\textbf{数据脱敏}：在输入AI前自动识别并替换敏感信息；\textbf{联邦学习}：多个医院可以在不共享用户病历的前提下协同训练诊断模型；\textbf{差分隐私}：在模型输出或数据查询结果中注入受控的随机噪声。


\subsubsection*{\textbf{•伦理风险案例}}
伦理风险案例：AI招聘的偏见问题。一家公司使用AI筛选简历，AI学习了历史招聘数据，但历史数据中管理层90\%是男性。结果AI学会了对男性候选人给予更高的评分。这不是AI故意歧视，而是它忠实地学习了训练数据中隐含的社会偏见。

这个问题的严重性在于，AI系统不仅继承了偏见，还会将其放大。相关研究显示，当使用未经过偏见处理的AI筛选简历时，女性候选人获得面试邀请的概率比同等资历的男性候选人低约30\%。更令人担忧的是，AI系统还能挖掘出人类不会注意到的代理歧视，比如模型发现某所大学的毕业生离职率较低，于是自动提高该大学候选人的评分，但客观上导致了对其他院校毕业生的不公平歧视。解决这一问题需要在数据收集阶段就确保多样性，在模型训练中引入公平性约束，并且在部署后持续监控。

\begin{funbox}
\textbf{Tay聊天机器人的失控风波。} 

2016年3月，微软在Twitter上推出了AI聊天机器人Tay，设计成一位活泼的青少年女性。微软的初衷是让Tay在与用户的互动中学习，变得越来越聪明。然而上线不到24小时，Tay就在部分恶意用户的调教下学会了发表种族歧视、性别歧视等不当言论。微软被迫紧急下线Tay。这个事件揭示了AI安全管控中一个深刻的教训：当AI从受控环境走向开放世界，缺乏权限边界和安全过滤的后果可能是灾难性的。Tay的失控不是因为它变坏了，而是因为它忠实地学习了环境中存在的偏见，这与今天智能体系统中的权限管理、数据筛选和责任归属问题一脉相承。
\end{funbox}



% 责任归属原则有三个典型场景。在医疗诊断中，AI提供初步筛查结果，医生做出最终诊断并承担法律责任，AI是辅助工具，医生是责任主体。在自动驾驶L2/L3中，系统同时控制加速、转向和制动，但驾驶员必须持续监控并随时准备接管，发生事故时责任在驾驶员。在AI辅助编程中，开发者审查AI生成的代码并确认后提交，对代码的正确性负责。

\subsection{人机协作与责任归属}

当AI犯错时，谁负责?这是智能体走向实际应用时必须回答的问题。

\subsubsection*{\textbf{•人机协作五个层次}}

L0(纯人工)：无AI参与。
L1(AI建议，人决策)：AI提供参考信息，人类做最终决定。
L2(AI决策，人可否决)：AI自主执行常规操作，但在关键节点提供否决权。
L3(AI执行，人监督)：AI全流程自主执行，人类监控异常情况。
L4(完全自主)：AI独立完成所有操作，人类仅在必要时介入。 如图\ref{fig:ch3_3_32}所示，人机协作从L0纯人工到L4完全自主分为五个层次，风险越高人类介入程度越深。

\begin{figure}[H]
\centering
\BookFigureImage{./images_chapter_01-03/3.32.png}
\caption{人机协作的五个层次}
\label{fig:ch3_3_32} 
\end{figure}

\vfill
\begin{knowledgebox}
人机协作五层次在实际场景中的应用对比：

\begin{formalTable}[H]
\caption{人机协作五层次对比}
\label{tab:ch3_07_collaboration}
\begin{tabular}{c c c c}
\toprule
\bfseries 层次 & \bfseries 描述 & \bfseries 适用场景 & \bfseries 责任主体 \\ \midrule
L0 纯人工 & 无AI参与 & 创意策划、伦理决策 & 人类 \\
L1 建议-决策 & AI建议，人决策 & 医疗诊断、投资建议 & 人类 \\
L2 决策-否决 & AI自主，人可否决 & 智能客服、邮件分类 & 人类(监督) \\
L3 执行-监督 & AI全流程，人监控 & 自动驾驶L3、工厂自动化 & 人类/AI(视情况) \\
L4 完全自主 & AI独立执行 & 太空探索、深海作业 & AI/系统设计方 \\ \bottomrule
\end{tabular}
\end{formalTable}

选择协作层次的核心原则：风险越高，人类介入越深。
\end{knowledgebox}
\vspace{1em}


\subsubsection*{\textbf{•AI智能客服安全管控案例}}
某电商平台部署了AI客服，可以自动处理退换货、查询订单等常规操作。但有一次用户输入“把我的收货地址改成XXX”，AI执行后发现那是一个从未配送过的偏远地区。系统检测到这次地址修改属于高风险操作，\textbf{自动触发人机回环}\footnote{人机回环(Human-in-the-Loop)：关键决策点由人类参与或确认的人机协作模式。}机制，转人工客服确认。人工客服联系用户后发现是账户被盗。正是这个改地址操作提前预警了安全风险。这个案例生动展示了L2层级(AI决策、人可否决)在实际应用中的价值：常规操作AI自主执行，但涉及敏感操作时保留人类否决权。


\subsubsection*{\textbf{•责任归属原则}}

在医疗诊断中，AI提供初步筛查结果，医生做出最终诊断并承担法律责任，AI是辅助工具，医生是责任主体。在自动驾驶L2/L3中，系统同时控制加速、转向和制动，但驾驶员必须持续监控并随时准备接管，发生事故时责任在驾驶员。在AI辅助编程中，开发者审查AI生成的代码并确认后提交，对代码的正确性负责。

人机回环：关键决策点必须有人类参与。AI的置信度低于阈值时，自动转给人类处理。例如银行信贷审批：(1) 用户提交贷款申请，(2) AI Agent自动评估，(3) 低风险申请自动批准，(4) 高风险申请自动拒绝，(5) 中等风险申请转人工审核。


\subsubsection*{\textbf{•可解释AI与信任建立}}
\textbf{建立信任：\textbf{可解释AI}\footnote{可解释AI(Explainable AI， XAI)：能向人类展示其决策依据和推理过程的AI系统。}} 人类需要理解AI的决策过程，才能信任它。AI诊断系统不仅给出结论，还通过热力图标注出影响判断的图像区域。AI投资建议不仅给出评级，还列出分析逻辑和每个判断依据。

% 本章从提示词的文法与基本规范起步，逐步深入到角色设定、上下文控制和思维链推理等核心技术，进而探讨了Skill模块如何将零散的成功经验固化为可复用的能力单元。在此基础上，我们从被动问答升级到智能体的主动规划执行闭环，理解了感知、任务分解、工具调用与反馈调整构成的自主系统。最后，本章走出了数字世界的边界，进入具身智能与物理闭环的领域，直面实时性、鲁棒性和安全性的全新挑战，并系统梳理了权限管控、数据隐私、人机协作与责任归属等关键议题。从学会与AI对话到理解AI的边界，本章完成了一条从工具使用到系统认知的完整学习路径。




\subsubsection*{小结}

具身智能让AI从数字世界走向物理世界，面临实时性、鲁棒性和安全性的全新挑战。权限分级模型、最小权限原则和沙箱机制是保障智能体安全的核心手段，而人机协作的五个层次（L0至L4）遵循风险越高人类介入越深的核心原则。数据隐私保护则依赖本地部署、数据脱敏、联邦学习和差分隐私等核心技术的综合运用。


\bigskip
\subsection*{课后习题}

\begin{enumerate}
\setlength{\itemsep}{6pt}
\item 具身智能与虚拟AI在核心能力上有哪些本质差异？

\item 什么是物理闭环？请以无人机稳定飞行为例说明其工作过程。
\item Sim-to-Real技术面临的核心挑战是什么？域随机化如何帮助解决这一问题？
\item 智能体权限分级模型中，四个层级分别对应什么权限范围？
\item 人机协作的五个层次（L0至L4）是如何划分的？风险越高时人类介入程度应如何变化？
\end{enumerate}

% \section*{课后习题}
% % \addcontentsline{toc}{section}{课后习题}

% \begin{enumerate}
% \item 简述提示词工程中词法和语法的含义和各自的要素，并各举一个正面和反面的例子。

% \item 在以下提示词中，指出使用了哪些提示词技巧，并评价其优缺点：
% \begin{userbox}
% 帮我写一份分析报告。
% \end{userbox}

% \item 比较思维链(CoT)、思维树(ToT)和反思链三种推理技术的适用场景和核心区别。

% \item 假设需要设计一个实验数据处理Skill，请按照Skill蒸馏的方法，列出回顾、归纳、泛化三个步骤中应该关注的内容。

% \item 请以撰写一篇课程论文文献综述为例，描述一个智能体从接收到任务到完成任务的全过程，包括感知、规划、工具调用、反馈调整等环节。

% \item 分析具身智能与虚拟AI在安全风险方面的主要差异，并说明为什么具身智能的隐私保护面临更大挑战。

% \item 在AI辅助医疗诊断的场景中，人机协作应属于哪个层次?请说明理由，并分析这种情况下责任的归属问题。
% \item 假设你正在设计一个智能学术助手系统，需要协调三个智能体(文献搜索Agent、数据分析Agent、论文写作Agent)协作完成一篇文献综述。请设计它们之间的通信机制，并说明当一个Agent执行失败时的容错策略。
% \end{enumerate}
% \end{itemize}

\section*{全章回顾}

本章围绕提示词工程与智能体系统展开，从人与AI的有效对话出发，沿着从提示词设计到智能体系统的递进路径，构建了一条从工具使用到系统认知的完整学习路径。

\textbf{第3.1节 提示词的文法与设计}介绍了提示词工程的核心技能。词法决定了动词引导输出类型、形容词控制风格深度、约束条件限定内容边界；语法通过总分总、递进和对比等结构组织信息。角色设定的本质是在AI的知识空间中划定注意力边界，有效的角色设定包含身份、经验年限、沟通风格、思考框架和约束条件五大要素。上下文窗口的管理通过背景注入、示例引导和格式说明三种策略实现。思维链通过将一次性映射分解为多个中间推理步骤来提升复杂问题的准确率，其三种变体（零样本思维链、思维树、反思链）各有侧重。结构化输出通过指定表格、JSON、Mermaid等格式，让AI的结果可直接被程序消费和使用。

\textbf{第3.2节 Skill模块的抽取复用}讲述了如何将成功的AI交互经验固化为可复用的能力单元。Skill是提示词、工具和流程三者的结合体。Skill蒸馏三步法（回顾→归纳→泛化）从日常AI交互中提炼通用模板。多个子Skill可串联为流水线完成复杂任务，参数化调用让同一Skill通过不同配置适配多种场景。Coze、Dify等主流平台提供了从封装到调用的完整工具链，Skill库的分类体系、版本管理和分享协作机制进一步提升了团队效率。

\textbf{第3.3节 智能体规划执行闭环}讲解了AI如何从被动问答升级为自主完成任务。智能体的感知阶段经过字面解析、意图推断和上下文融合三层处理，信息不足时主动追问。任务规划将大目标拆解为有依赖关系的子任务，采用深度优先、广度优先或动态规划策略。执行阶段调用搜索、计算、API、文件等外部工具来完成具体操作。反馈调整通过自我反思和外部信号持续优化输出质量。多智能体系统采用主从、议会和市场三种组织模式来协调分工与解决冲突。

\textbf{第3.4节 具身智能与安全边界}探讨了AI从数字世界走向物理世界面临的全新挑战。具身智能与虚拟AI的核心差异在于实时性、鲁棒性和安全性三大方面。物理闭环（感知→规划→控制）每秒钟循环数百次，任何延迟都可能导致物理损坏。Sim-to-Real技术通过域随机化缓解仿真与真实的差距。智能体的权限分级模型（只读级、执行级、管理级、超级级）、最小权限原则、沙箱机制和人机回环是保障安全的核心手段。人机协作的五个层次（L0至L4）遵循风险越高人类介入越深的原则，而隐私保护则依赖本地部署、数据脱敏、联邦学习和差分隐私等核心技术。

本章四节内容层层递进：从掌握提示词的设计文法，到将成功经验固化为可复用的Skill模块，再到理解智能体的自主执行闭环，最后直面具身智能的安全边界与责任归属。这构成了从学会与AI对话到理解AI系统的完整知识体系，为读者在日常工作中高效、安全地使用AI奠定了坚实基础。


